\documentclass[11pt]{article}
\usepackage[spanish,es-nodecimaldot,es-noshorthands]{babel}
\usepackage[utf8]{inputenc}
\usepackage[T1]{fontenc}
\usepackage{lmodern,microtype}
\usepackage{booktabs,longtable,tabularx,amsmath,amssymb,float,xurl,placeins,xcolor}
\usepackage{pgfplots}
\usetikzlibrary{arrows.meta,positioning,shapes.geometric,fit,backgrounds,calc,babel}
\usepackage[font=small,labelfont=bf]{caption}
\usepackage[hidelinks]{hyperref}
\usepackage[a4paper,margin=2.3cm]{geometry}
\pgfplotsset{compat=1.18}
\setlength{\emergencystretch}{2em}
\renewcommand{\arraystretch}{1.12}
\setlength{\parskip}{0.4em}
\setlength{\parindent}{0pt}


\title{\textbf{Laboratorio 2}\\[2pt] Métodos supervisados: clasificación\\[2pt]\large KNN, árboles de decisión, Random Forest, XGBoost y AdaBoost}
\author{Dylan Elizondo Alvarado (504610652)\\Joseph Garcia Montero (402680547)\\Rodrigo Ureña Castillo (118910482)\\Mariana Ureña Padilla (119180324)}
\date{EIF420O -- Inteligencia Artificial, I Ciclo 2026\\Universidad Nacional de Costa Rica\\Profesor: Dr. Juan Murillo-Morera\\15 de octubre de 2026}


\begin{document}

\maketitle

\begin{abstract}

Se compararon cinco familias de clasificadores (KNN, árbol de decisión, Random
Forest, XGBoost y AdaBoost) en dos problemas: potabilidad del agua (2011 muestras,
9 variables fisicoquímicas) y diabetes (390 pacientes, 17 variables clínicas). Para
cada familia se evaluó la configuración estándar y tres variantes que cambian un
único hiperparámetro, con validación cruzada estratificada de 5 pliegues sobre el
75\,\% de los datos, y se probó además la selección de atributos por información
mutua. El modelo elegido en cada familia se evaluó una sola vez en un 25\,\% de
prueba reservado. En Potabilidad, Random Forest con 500 árboles obtuvo el mejor
resultado (F1 macro 0.659, exactitud 70.2\,\%, AUC 0.727), pero ningún modelo
detectó más de la mitad de las muestras potables: la señal de los datos es débil y
no lineal. En Diabetes, todos los modelos basados en árboles alcanzaron F1 macro
entre 0.76 y 0.78 y AUC cercanas a 0.9, dominados por la glucosa estabilizada; los
siete diabéticos que ningún modelo detecta tienen glucosa en rango normal, lo que
muestra un límite de la información disponible y no de los algoritmos. Todo el
trabajo se integró en la clase \texttt{Clasificacion} del framework del curso y es
reproducible con un solo comando por dataset.


\end{abstract}

\tableofcontents

\clearpage

\section{Introducción y objetivos}

El objetivo general del laboratorio es realizar experimentos de clasificación
supervisada con los conjuntos de datos del curso, Potabilidad y Diabetes, usando el
framework de Python construido durante el semestre. Los objetivos específicos son:

\begin{enumerate}
\item Explicar el funcionamiento de KNN, árboles de decisión, Random Forest,
XGBoost y AdaBoost (Sección~2).
\item Definir modelos que combinen preprocesamiento, selección de atributos y cada
algoritmo; para cada uno, comparar la configuración estándar con variaciones,
seleccionar el mejor modelo y analizarlo (Secciones~3 a~5).
\end{enumerate}

Para que las comparaciones sean justas y repetibles se fijaron tres reglas. Primero,
cada variante cambia un solo hiperparámetro respecto del estándar, de modo que
cualquier diferencia se pueda atribuir a ese cambio. Segundo, la elección se hace
con validación cruzada sobre el conjunto de entrenamiento, y el conjunto de prueba
se usa una única vez al final. Tercero, todo el preprocesamiento se aprende sólo
con datos de entrenamiento, para evitar fuga de información.

El documento se organiza así: la Sección~2 presenta los fundamentos teóricos y un
mapa conceptual; la Sección~3, la metodología y su implementación en el framework;
las Secciones~4 y~5, los resultados y el análisis de cada dataset; la Sección~6
compara ambos casos y discute limitaciones, y la Sección~7 resume las conclusiones.
El anexo contiene los comandos para reproducir todo.


\section{Fundamentos teóricos}

\subsection{Aprendizaje supervisado y generalización}
En clasificación, cada observación contiene predictoras $x_i\in\mathbb{R}^d$ y una
etiqueta discreta $y_i$. El modelo aprende una regla $\hat y=f(x)$ a partir de
datos de entrenamiento, y su utilidad se mide en observaciones que no participaron
en ese aprendizaje \cite{hastie}. Los \emph{hiperparámetros} (por ejemplo, el número
de vecinos o la profundidad de un árbol) controlan el procedimiento de ajuste; los
\emph{parámetros} se estiman con los datos. Más complejidad reduce el error de
entrenamiento, pero puede aumentar la varianza ante muestras nuevas: el modelo
memoriza ruido (sobreajuste). Si el modelo es demasiado simple, no captura la
estructura (subajuste). Por eso hay que separar ajuste, selección y evaluación final
\cite{sklearnleakage,sklearnmetrics}.

\subsection{KNN: vecinos más cercanos}
KNN no construye un modelo explícito: guarda las observaciones de entrenamiento y
clasifica un punto nuevo según las etiquetas de sus $k$ vecinos más próximos
\cite{coverhart,sklearnknn}. Con la distancia de Minkowski,
\[
 d_p(x,z)=\Bigl(\sum_{j=1}^{d}|x_j-z_j|^p\Bigr)^{1/p},
 \qquad
 \hat y=\arg\max_c\sum_{i\in N_k(x)}w_i\,\mathbf{1}(y_i=c),
\]
donde $p=2$ es la distancia euclidiana, $p=1$ la Manhattan, y los pesos $w_i$ son
uniformes o inversamente proporcionales a la distancia. Un $k$ pequeño conserva
detalles locales y es sensible al ruido; uno grande suaviza la frontera y puede
diluir una clase poco frecuente, porque la mayoría de los vecinos pertenecerá a la
clase dominante. Como todas las variables entran en la distancia con el mismo peso,
el escalado es obligatorio y los atributos irrelevantes empeoran el método: con
muchas dimensiones las distancias tienden a parecerse entre sí.

\subsection{Árbol de decisión (DT)}
Un árbol divide recursivamente el espacio de atributos con reglas de la forma
$x_j\le t$. Cada hoja es una región y predice la clase más frecuente de las
observaciones de train que caen en ella. La implementación utilizada,
\texttt{DecisionTreeClassifier}, sigue el algoritmo CART \cite{sklearntree}. Para las
proporciones $p_c$ de cada clase en un nodo, las impurezas habituales son
\[
 G=1-\sum_c p_c^2 \quad\text{(Gini)}, \qquad H=-\sum_c p_c\log p_c \quad\text{(entropía)},
\]
y en cada nodo se elige la división que más reduce la impureza ponderada de los
hijos. Sin restricciones, el árbol crece hasta que cada hoja es pura y memoriza el
entrenamiento. La profundidad máxima y el mínimo de observaciones por hoja limitan
esa complejidad. Sus reglas se pueden leer directamente, pero cambios pequeños en la
muestra producen árboles distintos: es un método de alta varianza.

\subsection{Random Forest (RF)}
Random Forest promedia muchos árboles profundos, cada uno entrenado con una muestra
\emph{bootstrap} de las filas y con un subconjunto aleatorio de atributos candidatos
en cada división \cite{breiman,sklearnensemble}. Cada árbol individual sobreajusta,
pero sus errores están poco correlacionados y el promedio reduce la varianza sin
aumentar mucho el sesgo. \texttt{n\_estimators} fija cuántos árboles se promedian
(más árboles estabilizan la predicción pero no cambian el sesgo), y
\texttt{max\_depth} o \texttt{min\_samples\_leaf} limitan cada árbol. Es habitual que
un bosque alcance exactitud casi perfecta en train sin que eso implique la misma
calidad en datos nuevos.

\subsection{XGBoost}
XGBoost construye un ensamble \emph{aditivo}: cada árbol nuevo corrige los errores
del conjunto anterior siguiendo el gradiente de una función de pérdida, con una
penalización explícita de complejidad \cite{chen,xgboostmodel}:
\[
 F_T(x)=\sum_{t=1}^{T}\eta f_t(x),\qquad
 \mathcal L=\sum_i\ell\bigl(y_i,F_T(x_i)\bigr)+\sum_t\Omega(f_t),
 \qquad \Omega(f)=\gamma L+\tfrac{\lambda}{2}\sum_{j=1}^{L}w_j^2,
\]
donde $L$ es el número de hojas, $w_j$ el valor de cada hoja y $\eta$ la tasa de
aprendizaje. Una $\eta$ menor hace que cada árbol aporte menos y actúa como
regularización; una profundidad menor limita las interacciones que cada árbol
puede representar. Se usa la biblioteca \texttt{xgboost.XGBClassifier}, no el
\texttt{GradientBoostingClassifier} de scikit-learn.

\subsection{AdaBoost}
AdaBoost combina en secuencia aprendices débiles, por defecto árboles de un solo
nivel (\emph{stumps}), y aumenta el peso de las observaciones mal clasificadas para
que el siguiente aprendiz se concentre en ellas \cite{freund,sklearnada}. En la
versión binaria, con error ponderado $\epsilon_t$,
\[
 \alpha_t=\tfrac12\log\frac{1-\epsilon_t}{\epsilon_t},\qquad
 D_{t+1}(i)\propto D_t(i)\,e^{-\alpha_t y_i h_t(x_i)},\qquad
 \hat y=\operatorname{sign}\Bigl(\sum_t\alpha_t h_t(x)\Bigr).
\]
El número de estimadores y la tasa de aprendizaje regulan el ensamble. Si cada stump
apenas supera el azar ($\epsilon_t\approx0.5$), su peso $\alpha_t$ es casi nulo y el
ensamble aprende muy poco; además, concentrar peso en casos difíciles puede
amplificar el ruido de las etiquetas.

\subsection{Evaluación: matriz de confusión y métricas}
La matriz de confusión cuenta $C_{ij}=\#\{y=i,\hat y=j\}$ (filas reales, columnas
predichas). Para cada clase $c$, tratada como positiva frente al resto,
\[
 P_c=\frac{TP_c}{TP_c+FP_c},\qquad
 R_c=\frac{TP_c}{TP_c+FN_c},\qquad
 F1_c=\frac{2P_cR_c}{P_c+R_c},\qquad
 \text{F1 macro}=\frac1{|C|}\sum_c F1_c .
\]
La exactitud $\sum_c C_{cc}/n$ puede ser alta aunque una clase poco frecuente no se
detecte nunca. El F1 macro da el mismo peso a cada clase y por eso es la métrica de
selección de este trabajo: los dos conjuntos están desbalanceados (60/40 y 85/15)
\cite{sklearnmetrics,hegarcia}.

La \textbf{curva ROC} grafica la tasa de verdaderos positivos contra la de falsos
positivos al mover el umbral de decisión sobre la probabilidad estimada. El área
bajo la curva (AUC) es la probabilidad de que el modelo asigne mayor puntuación a un
positivo que a un negativo elegidos al azar: 0.5 equivale al azar y 1 a una
separación perfecta. No depende del umbral 0.5 que se usa para producir la matriz
de confusión \cite{fawcett}.

\subsection{Validación cruzada estratificada}
Una única partición de validación es muy sensible a qué filas caen en ella,
sobre todo con pocos casos de la clase minoritaria. En la validación cruzada de $K$
pliegues, el entrenamiento se divide en $K$ partes; cada configuración se entrena $K$
veces con $K-1$ partes y se evalúa en la restante, y se reporta la media y la
desviación de la métrica \cite{kohavi,sklearncv}. La versión \emph{estratificada}
conserva la proporción de clases en cada pliegue. Todo preprocesamiento (imputación,
codificación, escalado y selección de atributos) se ajusta dentro de cada pliegue,
para que el pliegue de evaluación nunca influya en el modelo. El conjunto de prueba
se separa antes y sólo se usa una vez, con el modelo ya elegido.

\subsection{Selección de atributos por información mutua}
La información mutua $I(X_j;Y)$ mide cuánto se reduce la incertidumbre sobre la clase
al conocer una variable; es cero si son independientes y, a diferencia de la
correlación, también detecta relaciones no lineales \cite{kraskov,ross}. Un método de
filtro como \texttt{SelectKBest} conserva las $k$ variables con mayor $I(X_j;Y)$
\cite{guyon}. Su limitación es que evalúa cada variable por separado: una variable que
sólo es útil en combinación con otra puede recibir una puntuación cercana a cero.

\subsection{Desbalance de clases y ponderación}
Con clases desbalanceadas, un algoritmo que minimiza el error total tiende a favorecer
la clase mayoritaria. Una corrección simple es ponderar cada observación de train con
$w_i=\frac{n}{|C|\,n_{y_i}}$, inversamente proporcional a la frecuencia de su clase,
de modo que ambas clases pesen lo mismo en el ajuste \cite{hegarcia}. Se aplica a DT, RF,
XGBoost y AdaBoost; KNN no admite pesos por observación.

\subsection{Importancia por permutación}
Para interpretar un modelo ya entrenado se permuta al azar una variable en el conjunto
de prueba y se mide cuánto empeora la métrica. Si la caída es grande, el modelo
depende de esa variable; si es nula, la ignora. Al ser independiente del algoritmo,
permite comparar KNN con los métodos basados en árboles \cite{breiman,sklearnperm}.
No demuestra causalidad.

\subsection{Mapa conceptual}
La Figura~\ref{fig:mapa} resume cómo se relacionan los conceptos anteriores dentro del
flujo experimental.

\begin{figure}[H]\centering
\begin{tikzpicture}[
  font=\footnotesize,
  caja/.style={draw, rounded corners=2pt, align=center, inner sep=4pt, anchor=north, fill=#1},
  caja/.default=white,
  flecha/.style={-{Stealth[length=2mm]}, thick, gray!70!black}]
\node[caja=gray!12, text width=13.2cm] (sup) at (0.5,0) {\textbf{Aprendizaje supervisado: clasificación}\\ datos etiquetados $(x_i,y_i)$ $\rightarrow$ regla $\hat y=f(x)$ que generalice a datos nuevos};
\node[caja=blue!8, text width=5.0cm] (inst) at (-3.6,-1.6) {\textbf{Basado en instancias}\\ KNN: vota con los $k$ vecinos;\\ depende de la distancia\\ y del escalado};
\node[caja=green!10, text width=5.6cm] (arb) at (3.7,-1.6) {\textbf{Basado en árboles}\\ DT: reglas $x_j\le t$ elegidas por impureza (Gini, entropía)};
\node[caja=yellow!15, text width=5.0cm] (pre) at (-3.6,-3.75) {\textbf{Preparación (sólo con train)}\\ imputación, one-hot, escalado,\\ selección por información mutua,\\ ponderación de clases};
\node[caja=green!18, text width=3.3cm] (bag) at (1.9,-3.75) {\textbf{Bagging}\\ RF: promedia árboles independientes\\ $\rightarrow$ menos varianza};
\node[caja=orange!15, text width=3.3cm] (boost) at (5.5,-3.75) {\textbf{Boosting}\\ secuencial, corrige errores: AdaBoost (pesos), XGBoost (gradiente)};
\node[caja=purple!10, text width=13.2cm] (val) at (0.5,-6.2) {\textbf{Selección del modelo}: estándar y variantes de un solo cambio $\rightarrow$ validación cruzada estratificada de 5 pliegues $\rightarrow$ ganador por familia con mayor F1 macro};
\node[caja=red!8, text width=13.2cm] (eval) at (0.5,-7.5) {\textbf{Evaluación final en test reservado}: matriz de confusión, precisión, recall y F1 por clase, ROC/AUC\\ \textbf{Interpretación}: importancia por permutación, reglas del árbol, sobreajuste (train vs.\ CV vs.\ test)};
\draw[flecha] (sup.south -| inst.north) -- (inst.north);
\draw[flecha] (sup.south -| arb.north) -- (arb.north);
\draw[flecha] (arb.south -| bag.north) -- (bag.north);
\draw[flecha] (arb.south -| boost.north) -- (boost.north);
\draw[flecha] (inst) -- (pre);
\draw[flecha] (pre.south) -- (pre.south |- val.north);
\draw[flecha] (bag.south) -- (bag.south |- val.north);
\draw[flecha] (boost.south) -- (boost.south |- val.north);
\draw[flecha] (val) -- (eval);
\end{tikzpicture}
\caption{Mapa conceptual: familias de clasificadores, preparación sin fuga de información, selección por validación cruzada y evaluación final.}
\label{fig:mapa}
\end{figure}


\clearpage

\section{Metodología}

\subsection{Integración en el framework}
Los algoritmos se implementaron en la jerarquía de clases del framework del curso:
\texttt{EDA} $\rightarrow$ \texttt{Supervisado} $\rightarrow$ \texttt{Clasificacion}
(módulo \texttt{framework\_ia/modelos/supervisado}). \texttt{Clasificacion} expone un
método por algoritmo (\texttt{knn}, \texttt{arbol\_decision}, \texttt{random\_forest},
\texttt{xgboost}, \texttt{adaboost}, con los alias \texttt{KNN}, \texttt{DT},
\texttt{RF}, \texttt{XGBoost}, \texttt{AdaBoost}) para entrenar una configuración, y el
método \texttt{experimentar} para el protocolo completo de este informe:

\begin{verbatim}
from framework_ia.modelos import Clasificacion
modelo = Clasificacion(dataframe=datos, target="diabetes", features=predictoras)
resultado = modelo.experimentar(estrategia="cv", folds=5, random_state=42,
                                incluir_categoricas=True)
resultado.tabla            # todas las configuraciones evaluadas
resultado.mejores["RF"]    # ganador de Random Forest evaluado en test
\end{verbatim}

El preprocesamiento vive en \texttt{PreprocesadorSupervisado}: imputación por mediana
(numéricas) y moda (categóricas), codificación one-hot, estandarización y, de forma
opcional, selección de los $k$ atributos con mayor información mutua. Es un
transformador de scikit-learn \cite{pedregosa} que se ajusta en cada pliegue sólo con
las filas de entrenamiento. La validación cruzada está en
\texttt{validacion\_cruzada.py}, y el script \texttt{scripts/ejecutar\_lab02.py} ejecuta
el mismo flujo desde la terminal y exporta toda la evidencia (tablas, matrices,
predicciones, curvas ROC, importancias y reglas del árbol). La interfaz Streamlit del
framework permite entrenar cada algoritmo de forma individual y, en la vista
\emph{Experimentos Lab02}, ejecutar este mismo protocolo con validación cruzada.

\subsection{Preparación de los datos}
\textbf{Potabilidad.} El archivo público de Water Quality \cite{datospotabilidad} tiene
3276 filas, pero \texttt{ph}, \texttt{Sulfate} y \texttt{Trihalomethanes} tienen 491,
781 y 162 valores ausentes. La versión usada en clase contiene 2011 filas, que son
exactamente los casos completos del archivo; por eso se trabajó con esos 2011 casos,
sin imputar. El objetivo \texttt{Potability} vale 1 para agua potable y 0 para no potable.

\textbf{Diabetes.} El archivo original \cite{datosdiabetes} tiene 403 pacientes y 19
columnas, sin una columna de clase. La clase se construyó con el criterio histórico del
propio conjunto: hemoglobina glicosilada \texttt{glyhb} $>7$ indica diabetes. Se
excluyeron las 13 filas sin \texttt{glyhb} (no se imputa la etiqueta), y se retiraron
\texttt{id} (identificador) y \texttt{glyhb} (define la clase: usarla sería fuga
directa). Quedan 390 pacientes y 17 predictoras: 14 numéricas y 3 categóricas
(\texttt{location}, \texttt{gender}, \texttt{frame}). Los valores ausentes de las
predictoras (por ejemplo, 252 en \texttt{bp.2s} y \texttt{bp.2d}) se imputan dentro de
cada ajuste.

Ambas preparaciones son deterministas, se ejecutan con un script cada una (Anexo) y
registran el SHA-256 de cada archivo.

\subsection{Protocolo experimental}
\begin{enumerate}
\item \textbf{Prueba reservada.} Se separa un 25\,\% estratificado (semilla 42) que no
interviene en ninguna decisión.
\item \textbf{Validación cruzada.} El 75\,\% restante se divide en 5 pliegues
estratificados. Cada configuración del catálogo (Tabla~\ref{tab:catalogo}) se entrena
5 veces y se registra la media y la desviación estándar del F1 macro en los pliegues
de validación. Se elige F1 macro porque ambos datasets están desbalanceados.
\item \textbf{Variantes de un solo cambio.} Cada familia tiene su configuración estándar
(valores predeterminados de la biblioteca) y tres variantes. Para DT, RF, XGBoost y
AdaBoost una de ellas pondera las clases de forma inversa a su frecuencia.
\item \textbf{Selección de atributos.} La mejor configuración de cada familia se vuelve a
evaluar con los mismos pliegues conservando sólo la mitad de los atributos con mayor
información mutua. Se adopta la versión con mejor F1 macro medio.
\item \textbf{Modelo final.} El ganador de cada familia se reajusta con todo el
entrenamiento y se evalúa una vez en prueba: matriz de confusión, métricas por clase,
exactitud, F1 macro y AUC.
\item \textbf{Interpretación.} Para el ganador se calcula la importancia por permutación
en prueba, las importancias internas de los árboles y, en DT, sus reglas. Estas
herramientas no se usan para seleccionar.
\end{enumerate}


\begin{table}[H]\centering\footnotesize
\caption{Catálogo de configuraciones. Cada variante cambia un único aspecto respecto del estándar, para poder atribuirle su efecto.}
\label{tab:catalogo}
\begin{tabular}{llp{0.27\linewidth}p{0.36\linewidth}}\toprule
Familia & Nombre & Parámetros solicitados & Qué cambia \\ \midrule
KNN & \texttt{estandar} & \texttt{--} & valores predeterminados de la biblioteca \\
KNN & \texttt{k\_\allowbreak{}15} & \texttt{n\_\allowbreak{}neighbors=15} & $k=15$ vecinos (estándar: 5) \\
KNN & \texttt{pesos\_\allowbreak{}distancia} & \texttt{weights=distance} & votos ponderados por $1/d$ (estándar: uniformes) \\
KNN & \texttt{manhattan} & \texttt{p=1} & distancia Manhattan, $p=1$ (estándar: euclidiana) \\
DT & \texttt{estandar} & \texttt{--} & valores predeterminados de la biblioteca \\
DT & \texttt{profundidad\_\allowbreak{}5} & \texttt{max\_\allowbreak{}depth=5} & profundidad máxima 5 (estándar: sin límite) \\
DT & \texttt{hoja\_\allowbreak{}minima\_\allowbreak{}10} & \texttt{min\_\allowbreak{}samples\_\allowbreak{}leaf=10} & mínimo 10 observaciones por hoja (estándar: 1) \\
DT & \texttt{clases\_\allowbreak{}balanceadas} & \texttt{balancear\_\allowbreak{}clases=True} & pesos por observación inversos a la frecuencia de su clase \\
RF & \texttt{estandar} & \texttt{--} & valores predeterminados de la biblioteca \\
RF & \texttt{500\_\allowbreak{}arboles} & \texttt{n\_\allowbreak{}estimators=500} & 500 árboles (estándar: 100) \\
RF & \texttt{hoja\_\allowbreak{}minima\_\allowbreak{}5} & \texttt{min\_\allowbreak{}samples\_\allowbreak{}leaf=5} & mínimo 5 observaciones por hoja (estándar: 1) \\
RF & \texttt{clases\_\allowbreak{}balanceadas} & \texttt{balancear\_\allowbreak{}clases=True} & pesos por observación inversos a la frecuencia de su clase \\
XGBoost & \texttt{estandar} & \texttt{--} & valores predeterminados de la biblioteca \\
XGBoost & \texttt{tasa\_\allowbreak{}0\_\allowbreak{}05} & \texttt{learning\_\allowbreak{}rate=0.05} & tasa de aprendizaje 0.05 (estándar: 0.3) \\
XGBoost & \texttt{profundidad\_\allowbreak{}3} & \texttt{max\_\allowbreak{}depth=3} & profundidad máxima 3 (estándar: 6) \\
XGBoost & \texttt{clases\_\allowbreak{}balanceadas} & \texttt{balancear\_\allowbreak{}clases=True} & pesos por observación inversos a la frecuencia de su clase \\
AdaBoost & \texttt{estandar} & \texttt{--} & valores predeterminados de la biblioteca \\
AdaBoost & \texttt{200\_\allowbreak{}estimadores} & \texttt{n\_\allowbreak{}estimators=200} & 200 estimadores (estándar: 50) \\
AdaBoost & \texttt{tasa\_\allowbreak{}0\_\allowbreak{}5} & \texttt{learning\_\allowbreak{}rate=0.5} & tasa de aprendizaje 0.5 (estándar: 1.0) \\
AdaBoost & \texttt{clases\_\allowbreak{}balanceadas} & \texttt{balancear\_\allowbreak{}clases=True} & pesos por observación inversos a la frecuencia de su clase \\
\bottomrule\end{tabular}
\end{table}

\begin{table}[H]\centering\footnotesize
\caption{Tamaños efectivos del protocolo (semilla 42).}
\begin{tabular}{lrrrcrrr}\toprule
Dataset & Filas & Train & Test & Validación por pliegue & Predictoras & Atributos transf. & $k$ sel. \\ \midrule
Potabilidad del agua & 2011 & 1508 & 503 & 302/302/302/301/301 & 9 & 9 & 5 \\
Diabetes & 390 & 292 & 98 & 59/59/58/58/58 & 17 & 21 & 11 \\
\bottomrule\end{tabular}
\end{table}

\FloatBarrier

\clearpage

\section{Resultados: Potabilidad del agua}

\subsection{Datos y partición}

Se analizan 2011 muestras de agua con nueve mediciones fisicoquímicas (\texttt{ph},
dureza, sólidos disueltos, cloraminas, sulfato, conductividad, carbono orgánico,
trihalometanos y turbidez). El 59.7\,\% de las muestras es no potable. Un rasgo
central de este conjunto es que las medias de cada variable son casi idénticas en las
dos clases (por ejemplo, \texttt{ph} 7.07 frente a 7.11 y sulfato 333.7 frente a
332.5\,mg/L): ninguna variable separa las clases por sí sola. Lo que cambia es la
dispersión y los extremos. El 73\,\% de las muestras con sulfato menor que 260\,mg/L
y el 62\,\% de las que superan 420\,mg/L son potables, frente al 35\,\% entre 300 y
340\,mg/L. Con el \texttt{ph} ocurre algo parecido: la proporción de potables es mayor
cerca de la neutralidad (43--45\,\% entre 6 y 8) que en los extremos (30\,\% con
\texttt{ph} $\le 5$). Son relaciones en forma de U, no lineales, y eso condiciona qué
algoritmos funcionan mejor.

\begin{table}[H]\centering\small
\caption{Procedencia de Potabilidad del agua.}
\begin{tabular}{lp{0.62\linewidth}}\toprule
Elemento & Valor \\ \midrule
Archivo analizado & \texttt{potabilidad.csv} \\
Original & \texttt{water\_\allowbreak{}potability.csv}, 3276 filas \\
Filas analizadas & 2011 \\
Objetivo & \texttt{Potability} \\
Predictoras & 9 \\
Fuente & \url{https://www.kaggle.com/datasets/adityakadiwal/water-potability} \\
SHA-256 analizado & {\footnotesize\texttt{d2569bece50025b9dbf21cfd720ecbed}\newline\texttt{cdaa838ca082db29ef853c5fb9da8e12}} \\
\bottomrule\end{tabular}
\end{table}

\begin{table}[H]\centering\small
\caption{Distribución de clases en Potabilidad del agua (partición estratificada 75/25, semilla 42).}
\begin{tabular}{lrrr}\toprule
Clase & Train & Test & Total \\ \midrule
No potable (0) & 900 & 300 & 1200 \\
Potable (1) & 608 & 203 & 811 \\
\textbf{Total} & 1508 & 503 & 2011 \\
\bottomrule\end{tabular}
\end{table}

Una referencia trivial que siempre predice la clase mayoritaria de train (No potable (0)) obtendría en test una exactitud de 59.6\% y un F1 macro de 0.374. Cualquier modelo útil debe superar con claridad ese F1 macro, no sólo la exactitud.

\subsection{K vecinos más cercanos (KNN)}

\begin{table}[H]\centering\footnotesize
\caption{K vecinos más cercanos (KNN) en Potabilidad del agua: validación cruzada estratificada de 5 pliegues sobre train. $\star$ = configuración elegida; $\Delta$ = diferencia de F1 macro frente al estándar.}
\begin{tabular}{>{\raggedright\arraybackslash}p{0.21\linewidth}>{\raggedright\arraybackslash}p{0.26\linewidth}rrrr}\toprule
Configuración & Cambio respecto del estándar & F1 macro CV & $\Delta$ & Exact. CV & Recall CV \\ \midrule
\texttt{estandar} & valores predeterminados de la biblioteca & 0.5892 $\pm$ 0.029 & -- & 62.6\% & 0.5901 \\
\texttt{k\_\allowbreak{}15} & $k=15$ vecinos (estándar: 5) & 0.6000 $\pm$ 0.030 & +0.011 & 65.6\% & 0.6066 \\
\texttt{pesos\_\allowbreak{}distancia} & votos ponderados por $1/d$ (estándar: uniformes) & 0.5877 $\pm$ 0.030 & -0.001 & 62.5\% & 0.5887 \\
\texttt{manhattan} $\star$ & distancia Manhattan, $p=1$ (estándar: euclidiana) & 0.6036 $\pm$ 0.030 & +0.014 & 64.1\% & 0.6042 \\
\texttt{manhattan+\allowbreak{}seleccion\_\allowbreak{}k5} & mejor anterior + 5 atributos por información mutua & 0.5364 $\pm$ 0.033 & -0.053 & 57.7\% & 0.5399 \\
\bottomrule\end{tabular}
\end{table}

\textbf{Modelo final.} Configuración \texttt{manhattan} con todos los atributos, reajustada con las 1508 filas de train. F1 macro: 0.7472 en entrenamiento, 0.6036 en validación cruzada y 0.5969 en test (exactitud 63.4\%, AUC 0.668).

\begin{table}[H]\footnotesize\begin{minipage}[t]{0.44\linewidth}\centering\captionof{table}{K vecinos más cercanos (KNN): matriz de confusión en test (filas: real; columnas: predicción).}\begin{tabular}{lrr}\toprule
Real / pred. & No potable (0) & Potable (1) \\ \midrule
No potable (0) & 236 & 64 \\
Potable (1) & 120 & 83 \\
\bottomrule\end{tabular}\end{minipage}\hfill\begin{minipage}[t]{0.52\linewidth}\centering\captionof{table}{K vecinos más cercanos (KNN): métricas por clase en test.}\begin{tabular}{lrrrr}\toprule
Clase & Soporte & Precisión & Recall & F1 \\ \midrule
No potable (0) & 300 & 0.6629 & 0.7867 & 0.7195 \\
Potable (1) & 203 & 0.5646 & 0.4089 & 0.4743 \\
\bottomrule\end{tabular}\end{minipage}\end{table}

\paragraph{Análisis.} Las cuatro configuraciones quedan dentro de una desviación estándar entre sí
(0.588--0.604, con desviaciones cercanas a 0.03), así que la elección de la distancia
Manhattan (+0.014 sobre el estándar) es una mejora pequeña y poco segura. Aumentar a
$k=15$ sube la exactitud de validación de 62.6\,\% a 65.7\,\% pero apenas el F1 macro:
promediar más vecinos suaviza la frontera hacia la clase mayoritaria. Ponderar por
distancia no cambia nada relevante. La selección de cinco atributos empeoró el F1 macro
a 0.536, por la razón que se discute en la síntesis. En prueba, KNN reconoce bien el
agua no potable (recall 0.787) pero sólo 83 de 203 muestras potables (recall 0.409).
La importancia por permutación reparte el efecto entre \texttt{ph}, sulfato y sólidos,
sin una variable dominante. Esto es coherente con el método: KNN usa las nueve
variables con el mismo peso, aunque la mayoría aporta poca información, y esas
dimensiones irrelevantes vuelven las distancias menos informativas. La diferencia entre
F1 de entrenamiento (0.747) y de prueba (0.597) muestra además que, con $k=5$, el
modelo depende mucho de los vecinos concretos.

\FloatBarrier

\subsection{Árbol de decisión (DT)}

\begin{table}[H]\centering\footnotesize
\caption{Árbol de decisión (DT) en Potabilidad del agua: validación cruzada estratificada de 5 pliegues sobre train. $\star$ = configuración elegida; $\Delta$ = diferencia de F1 macro frente al estándar.}
\begin{tabular}{>{\raggedright\arraybackslash}p{0.21\linewidth}>{\raggedright\arraybackslash}p{0.26\linewidth}rrrr}\toprule
Configuración & Cambio respecto del estándar & F1 macro CV & $\Delta$ & Exact. CV & Recall CV \\ \midrule
\texttt{estandar} & valores predeterminados de la biblioteca & 0.5654 $\pm$ 0.048 & -- & 57.5\% & 0.5681 \\
\texttt{profundidad\_\allowbreak{}5} & profundidad máxima 5 (estándar: sin límite) & 0.5346 $\pm$ 0.071 & -0.031 & 62.5\% & 0.5667 \\
\texttt{hoja\_\allowbreak{}minima\_\allowbreak{}10} $\star$ & mínimo 10 observaciones por hoja (estándar: 1) & 0.5908 $\pm$ 0.047 & +0.025 & 61.0\% & 0.5906 \\
\texttt{clases\_\allowbreak{}balanceadas} & pesos por observación inversos a la frecuencia de su clase & 0.5633 $\pm$ 0.049 & -0.002 & 58.4\% & 0.5633 \\
\texttt{hoja\_\allowbreak{}minima\_\allowbreak{}10+\allowbreak{}seleccion\_\allowbreak{}k5} & mejor anterior + 5 atributos por información mutua & 0.5487 $\pm$ 0.050 & -0.017 & 57.6\% & 0.5499 \\
\bottomrule\end{tabular}
\end{table}

\textbf{Modelo final.} Configuración \texttt{hoja\_\allowbreak{}minima\_\allowbreak{}10} con todos los atributos, reajustada con las 1508 filas de train. F1 macro: 0.7935 en entrenamiento, 0.5908 en validación cruzada y 0.5777 en test (exactitud 61.2\%, AUC 0.657).

\begin{table}[H]\footnotesize\begin{minipage}[t]{0.44\linewidth}\centering\captionof{table}{Árbol de decisión (DT): matriz de confusión en test (filas: real; columnas: predicción).}\begin{tabular}{lrr}\toprule
Real / pred. & No potable (0) & Potable (1) \\ \midrule
No potable (0) & 226 & 74 \\
Potable (1) & 121 & 82 \\
\bottomrule\end{tabular}\end{minipage}\hfill\begin{minipage}[t]{0.52\linewidth}\centering\captionof{table}{Árbol de decisión (DT): métricas por clase en test.}\begin{tabular}{lrrrr}\toprule
Clase & Soporte & Precisión & Recall & F1 \\ \midrule
No potable (0) & 300 & 0.6513 & 0.7533 & 0.6986 \\
Potable (1) & 203 & 0.5256 & 0.4039 & 0.4568 \\
\bottomrule\end{tabular}\end{minipage}\end{table}

\begin{figure}[H]\centering\begin{minipage}{0.92\linewidth}\scriptsize\begin{verbatim}
|--- Sulfate <= 261.02
|   |--- Solids <= 18346.62
|   |   |--- clase: 0 [n=13]
|   |--- Solids > 18346.62
|   |   |--- ph <= 6.23
|   |   |   |--- clase: 0 [n=10]
|   |   |--- ph > 6.23
|   |   |   |--- clase: 1 [n=39] (rama truncada)
|--- Sulfate > 261.02
|   |--- Sulfate <= 387.27
|   |   |--- ph <= 6.27
|   |   |   |--- clase: 0 [n=374] (rama truncada)
|   |   |--- ph > 6.27
|   |   |   |--- clase: 0 [n=939] (rama truncada)
|   |--- Sulfate > 387.27
|   |   |--- ph <= 7.91
|   |   |   |--- clase: 1 [n=93] (rama truncada)
|   |   |--- ph > 7.91
|   |   |   |--- clase: 0 [n=40] (rama truncada)
\end{verbatim}\end{minipage}\caption{Primeros tres niveles del árbol final (profundidad total 14, 107 hojas). Los umbrales están en las unidades originales: la estandarización se invirtió con la media y desviación de train. $n$ es la cantidad de observaciones de train en el nodo.}\end{figure}

\paragraph{Análisis.} El árbol sin restricciones (estándar) obtiene F1 macro 0.565 con una exactitud de
apenas 57.5\,\%, menor que la de predecir siempre ``no potable'': crece hasta
memorizar el entrenamiento. Limitar la profundidad a 5 sube la exactitud (62.5\,\%)
pero baja el F1 macro (0.535), porque un árbol tan corto termina prediciendo casi
siempre la clase mayoritaria. La mejor opción es exigir al menos 10 observaciones por
hoja (0.591, +0.025): regulariza sin impedir que el árbol aprenda las reglas en U. El
árbol final sigue siendo grande (profundidad 14, 107 hojas) y el F1 cae de 0.794 en
entrenamiento a 0.578 en prueba, lo que confirma su alta varianza. Sus primeras reglas
son interpretables y coinciden con lo observado en los datos: con sulfato
$\le 261$\,mg/L, sólidos mayores que 18\,347\,mg/L y \texttt{ph} $>6.23$ predice potable; con sulfato entre 261 y
387\,mg/L predice no potable, y por encima de 387\,mg/L vuelve a predecir potable
cuando \texttt{ph} $\le 7.91$. El árbol representa la forma en U combinando dos cortes
sobre la misma variable y una interacción con el \texttt{ph}.

\FloatBarrier

\subsection{Random Forest (RF)}

\begin{table}[H]\centering\footnotesize
\caption{Random Forest (RF) en Potabilidad del agua: validación cruzada estratificada de 5 pliegues sobre train. $\star$ = configuración elegida; $\Delta$ = diferencia de F1 macro frente al estándar.}
\begin{tabular}{>{\raggedright\arraybackslash}p{0.21\linewidth}>{\raggedright\arraybackslash}p{0.26\linewidth}rrrr}\toprule
Configuración & Cambio respecto del estándar & F1 macro CV & $\Delta$ & Exact. CV & Recall CV \\ \midrule
\texttt{estandar} & valores predeterminados de la biblioteca & 0.6287 $\pm$ 0.031 & -- & 67.7\% & 0.6310 \\
\texttt{500\_\allowbreak{}arboles} $\star$ & 500 árboles (estándar: 100) & 0.6435 $\pm$ 0.022 & +0.015 & 69.0\% & 0.6442 \\
\texttt{hoja\_\allowbreak{}minima\_\allowbreak{}5} & mínimo 5 observaciones por hoja (estándar: 1) & 0.6151 $\pm$ 0.022 & -0.014 & 67.2\% & 0.6209 \\
\texttt{clases\_\allowbreak{}balanceadas} & pesos por observación inversos a la frecuencia de su clase & 0.6278 $\pm$ 0.045 & -0.001 & 65.5\% & 0.6262 \\
\texttt{500\_\allowbreak{}arboles+\allowbreak{}seleccion\_\allowbreak{}k5} & mejor anterior + 5 atributos por información mutua & 0.5436 $\pm$ 0.039 & -0.085 & 60.1\% & 0.5538 \\
\bottomrule\end{tabular}
\end{table}

\textbf{Modelo final.} Configuración \texttt{500\_\allowbreak{}arboles} con todos los atributos, reajustada con las 1508 filas de train. F1 macro: 1.0000 en entrenamiento, 0.6435 en validación cruzada y 0.6585 en test (exactitud 70.2\%, AUC 0.727).

\begin{table}[H]\footnotesize\begin{minipage}[t]{0.44\linewidth}\centering\captionof{table}{Random Forest (RF): matriz de confusión en test (filas: real; columnas: predicción).}\begin{tabular}{lrr}\toprule
Real / pred. & No potable (0) & Potable (1) \\ \midrule
No potable (0) & 266 & 34 \\
Potable (1) & 116 & 87 \\
\bottomrule\end{tabular}\end{minipage}\hfill\begin{minipage}[t]{0.52\linewidth}\centering\captionof{table}{Random Forest (RF): métricas por clase en test.}\begin{tabular}{lrrrr}\toprule
Clase & Soporte & Precisión & Recall & F1 \\ \midrule
No potable (0) & 300 & 0.6963 & 0.8867 & 0.7801 \\
Potable (1) & 203 & 0.7190 & 0.4286 & 0.5370 \\
\bottomrule\end{tabular}\end{minipage}\end{table}

\paragraph{Análisis.} Random Forest obtiene el mejor resultado de este dataset. El estándar (100 árboles)
alcanza 0.629 y pasar a 500 árboles sube el F1 macro a 0.644 y además baja la
desviación entre pliegues de 0.031 a 0.022: con más árboles el promedio es más estable,
que es exactamente lo que predice la teoría del bagging. Exigir 5 observaciones por
hoja empeora (0.615), porque limita los árboles profundos que el bosque necesita para
capturar las interacciones. Ponderar las clases no ayudó. En prueba logra exactitud de
70.2\,\% (contra 59.6\,\% de la referencia), F1 macro 0.659 y AUC 0.727. Es el modelo
más preciso al predecir potable (0.719: de cada 10 muestras que marca como potables,
7 lo son), pero sigue dejando pasar 116 de las 203 muestras potables (recall 0.429).
El F1 de entrenamiento es 1.000: cada árbol memoriza su muestra, y aun así el promedio
generaliza mejor que cualquier modelo individual. La importancia por permutación
(Figura~\ref{fig:perm-potabilidad}) identifica el sulfato (caída de 0.131 en F1) y el
\texttt{ph} (0.096) como las variables decisivas, seguidas de cloraminas, dureza y
sólidos.

\FloatBarrier

\subsection{XGBoost}

\begin{table}[H]\centering\footnotesize
\caption{XGBoost en Potabilidad del agua: validación cruzada estratificada de 5 pliegues sobre train. $\star$ = configuración elegida; $\Delta$ = diferencia de F1 macro frente al estándar.}
\begin{tabular}{>{\raggedright\arraybackslash}p{0.21\linewidth}>{\raggedright\arraybackslash}p{0.26\linewidth}rrrr}\toprule
Configuración & Cambio respecto del estándar & F1 macro CV & $\Delta$ & Exact. CV & Recall CV \\ \midrule
\texttt{estandar} & valores predeterminados de la biblioteca & 0.6220 $\pm$ 0.025 & -- & 65.1\% & 0.6210 \\
\texttt{tasa\_\allowbreak{}0\_\allowbreak{}05} & tasa de aprendizaje 0.05 (estándar: 0.3) & 0.6172 $\pm$ 0.024 & -0.005 & 66.3\% & 0.6196 \\
\texttt{profundidad\_\allowbreak{}3} & profundidad máxima 3 (estándar: 6) & 0.6208 $\pm$ 0.034 & -0.001 & 65.1\% & 0.6197 \\
\texttt{clases\_\allowbreak{}balanceadas} $\star$ & pesos por observación inversos a la frecuencia de su clase & 0.6392 $\pm$ 0.044 & +0.017 & 65.8\% & 0.6375 \\
\texttt{clases\_\allowbreak{}balanceadas+\allowbreak{}seleccion\_\allowbreak{}k5} & mejor anterior + 5 atributos por información mutua & 0.5439 $\pm$ 0.012 & -0.078 & 56.6\% & 0.5441 \\
\bottomrule\end{tabular}
\end{table}

\textbf{Modelo final.} Configuración \texttt{clases\_\allowbreak{}balanceadas} con todos los atributos, reajustada con las 1508 filas de train. F1 macro: 1.0000 en entrenamiento, 0.6392 en validación cruzada y 0.5896 en test (exactitud 61.0\%, AUC 0.671).

\begin{table}[H]\footnotesize\begin{minipage}[t]{0.44\linewidth}\centering\captionof{table}{XGBoost: matriz de confusión en test (filas: real; columnas: predicción).}\begin{tabular}{lrr}\toprule
Real / pred. & No potable (0) & Potable (1) \\ \midrule
No potable (0) & 210 & 90 \\
Potable (1) & 106 & 97 \\
\bottomrule\end{tabular}\end{minipage}\hfill\begin{minipage}[t]{0.52\linewidth}\centering\captionof{table}{XGBoost: métricas por clase en test.}\begin{tabular}{lrrrr}\toprule
Clase & Soporte & Precisión & Recall & F1 \\ \midrule
No potable (0) & 300 & 0.6646 & 0.7000 & 0.6818 \\
Potable (1) & 203 & 0.5187 & 0.4778 & 0.4974 \\
\bottomrule\end{tabular}\end{minipage}\end{table}

\paragraph{Análisis.} Las variantes de tasa de aprendizaje (0.05) y profundidad (3) quedan prácticamente
iguales al estándar (0.617--0.622). La única mejora es ponderar las clases (0.639,
+0.017), que obliga al modelo a prestar más atención a las muestras potables. En
prueba esto se ve como el segundo mayor recall de potables (0.478, 97 de 203), a costa
de más falsos positivos (90 frente a 34 de RF) y una exactitud de 61.0\,\%. XGBoost es
el modelo que más cae de validación (0.639) a prueba (0.590), y su F1 de entrenamiento
es 1.000: con la tasa 0.3 y profundidad 6 predeterminadas, los 100 árboles alcanzan a
memorizar el entrenamiento. En este dataset ruidoso, el boosting por gradiente termina
ajustando ruido, mientras que el promedio de RF lo atenúa.

\FloatBarrier

\subsection{AdaBoost}

\begin{table}[H]\centering\footnotesize
\caption{AdaBoost en Potabilidad del agua: validación cruzada estratificada de 5 pliegues sobre train. $\star$ = configuración elegida; $\Delta$ = diferencia de F1 macro frente al estándar.}
\begin{tabular}{>{\raggedright\arraybackslash}p{0.21\linewidth}>{\raggedright\arraybackslash}p{0.26\linewidth}rrrr}\toprule
Configuración & Cambio respecto del estándar & F1 macro CV & $\Delta$ & Exact. CV & Recall CV \\ \midrule
\texttt{estandar} & valores predeterminados de la biblioteca & 0.4714 $\pm$ 0.018 & -- & 59.7\% & 0.5224 \\
\texttt{200\_\allowbreak{}estimadores} & 200 estimadores (estándar: 50) & 0.5014 $\pm$ 0.017 & +0.030 & 59.9\% & 0.5337 \\
\texttt{tasa\_\allowbreak{}0\_\allowbreak{}5} & tasa de aprendizaje 0.5 (estándar: 1.0) & 0.4501 $\pm$ 0.016 & -0.021 & 61.5\% & 0.5286 \\
\texttt{clases\_\allowbreak{}balanceadas} $\star$ & pesos por observación inversos a la frecuencia de su clase & 0.5404 $\pm$ 0.039 & +0.069 & 55.2\% & 0.5456 \\
\texttt{clases\_\allowbreak{}balanceadas+\allowbreak{}seleccion\_\allowbreak{}k5} & mejor anterior + 5 atributos por información mutua & 0.5281 $\pm$ 0.011 & +0.057 & 54.4\% & 0.5314 \\
\bottomrule\end{tabular}
\end{table}

\textbf{Modelo final.} Configuración \texttt{clases\_\allowbreak{}balanceadas} con todos los atributos, reajustada con las 1508 filas de train. F1 macro: 0.5747 en entrenamiento, 0.5404 en validación cruzada y 0.5822 en test (exactitud 59.8\%, AUC 0.615).

\begin{table}[H]\footnotesize\begin{minipage}[t]{0.44\linewidth}\centering\captionof{table}{AdaBoost: matriz de confusión en test (filas: real; columnas: predicción).}\begin{tabular}{lrr}\toprule
Real / pred. & No potable (0) & Potable (1) \\ \midrule
No potable (0) & 200 & 100 \\
Potable (1) & 102 & 101 \\
\bottomrule\end{tabular}\end{minipage}\hfill\begin{minipage}[t]{0.52\linewidth}\centering\captionof{table}{AdaBoost: métricas por clase en test.}\begin{tabular}{lrrrr}\toprule
Clase & Soporte & Precisión & Recall & F1 \\ \midrule
No potable (0) & 300 & 0.6623 & 0.6667 & 0.6645 \\
Potable (1) & 203 & 0.5025 & 0.4975 & 0.5000 \\
\bottomrule\end{tabular}\end{minipage}\end{table}

\paragraph{Análisis.} AdaBoost estándar tiene el peor desempeño del estudio: F1 macro 0.471 con exactitud
59.7\,\%, igual a la proporción de la clase mayoritaria. Es decir, casi siempre
predice ``no potable''. Bajar la tasa a 0.5 lo empeora (0.450) y duplicar los
estimadores mejora poco (0.501). La razón está en sus aprendices: cada stump es un
único corte sobre una variable, y como las medias de ambas clases son casi iguales,
ningún corte supera por mucho el azar. En ese caso el peso $\alpha_t$ de cada stump es
pequeño y el ensamble no logra representar la interacción entre sulfato
y \texttt{ph}. La ponderación de clases es la variante con mayor efecto de todo el
dataset (+0.069, F1 0.540): al igualar el peso de ambas clases, los stumps dejan de
favorecer la mayoritaria. En prueba alcanza el recall más alto de potables (0.498),
pero con precisión 0.502 y AUC 0.615, la más baja. A diferencia de RF y XGBoost no
sobreajusta (F1 0.575 en entrenamiento y 0.582 en prueba): su problema es de sesgo
(subajuste), no de varianza. Casi toda su capacidad se apoya en el sulfato.

\FloatBarrier

\subsection{Comparación de los cinco modelos elegidos}

\begin{table}[H]\centering\footnotesize
\caption{Potabilidad del agua: resumen de los ganadores. La selección usó sólo la columna F1 CV; el mejor algoritmo por validación es RF.}
\begin{tabular}{llrrrrrr}\toprule
Familia & Configuración & F1 train & F1 CV & F1 test & Exactitud test & Recall Potable (1) & AUC \\ \midrule
KNN & \texttt{manhattan} & 0.7472 & 0.6036 & 0.5969 & 63.4\% & 0.4089 & 0.668 \\
DT & \texttt{hoja\_\allowbreak{}minima\_\allowbreak{}10} & 0.7935 & 0.5908 & 0.5777 & 61.2\% & 0.4039 & 0.657 \\
RF & \texttt{500\_\allowbreak{}arboles} & 1.0000 & 0.6435 & 0.6585 & 70.2\% & 0.4286 & 0.727 \\
XGBoost & \texttt{clases\_\allowbreak{}balanceadas} & 1.0000 & 0.6392 & 0.5896 & 61.0\% & 0.4778 & 0.671 \\
AdaBoost & \texttt{clases\_\allowbreak{}balanceadas} & 0.5747 & 0.5404 & 0.5822 & 59.8\% & 0.4975 & 0.615 \\
\bottomrule\end{tabular}
\end{table}

\begin{figure}[H]\centering
\begin{tikzpicture}
\begin{axis}[ybar, bar width=8pt, width=\linewidth, height=6.2cm, ymin=0, ymax=1.08,
 ylabel={F1 macro}, symbolic x coords={KNN,DT,RF,XGBoost,AdaBoost}, xtick=data,
 enlarge x limits=0.14, ymajorgrids, legend style={at={(0.5,-0.16)},anchor=north,legend columns=4,font=\small},
 tick label style={font=\small}]
\addplot[fill=gray!35, draw=gray!70!black] coordinates {(KNN,0.7472) (DT,0.7935) (RF,1.0000) (XGBoost,1.0000) (AdaBoost,0.5747)};
\addplot[fill=blue!55, draw=blue!70!black, error bars/.cd, y dir=both, y explicit] coordinates {(KNN,0.6036) +- (0,0.0303) (DT,0.5908) +- (0,0.0473) (RF,0.6435) +- (0,0.0218) (XGBoost,0.6392) +- (0,0.0441) (AdaBoost,0.5404) +- (0,0.0389)};
\addplot[fill=orange!70, draw=orange!80!black] coordinates {(KNN,0.5969) (DT,0.5777) (RF,0.6585) (XGBoost,0.5896) (AdaBoost,0.5822)};
\draw[dashed, red!70!black, thick] ({rel axis cs:0,0}|-{axis cs:KNN,0.3736}) -- ({rel axis cs:1,0}|-{axis cs:KNN,0.3736});
\legend{Entrenamiento, Validación cruzada (media $\pm$ d.e.), Test}
\end{axis}\end{tikzpicture}
\caption{Potabilidad del agua: F1 macro del modelo elegido en cada familia. La línea roja es la referencia que siempre predice la clase mayoritaria (F1 macro 0.374). La distancia entre la barra gris y la azul indica sobreajuste.}
\label{fig:resumen-potabilidad}
\end{figure}

\begin{figure}[H]\centering
\begin{tikzpicture}
\begin{axis}[width=0.72\linewidth, height=0.62\linewidth, xmin=0, xmax=1, ymin=0, ymax=1,
 xlabel={Tasa de falsos positivos}, ylabel={Tasa de verdaderos positivos}, grid=major,
 legend style={at={(0.98,0.02)}, anchor=south east, font=\footnotesize}]
\addplot[dashed, gray, forget plot] coordinates {(0,0) (1,1)};
\addplot[thick, blue!70!black] coordinates {(0.0000,0.0000) (0.0167,0.0197) (0.0600,0.1872) (0.2133,0.4089) (0.5200,0.7734) (0.8300,0.9655) (1.0000,1.0000)};
\addlegendentry{KNN (AUC 0.668)}
\addplot[thick, orange!85!black] coordinates {(0.0000,0.0000) (0.0433,0.1379) (0.0433,0.1626) (0.0500,0.1872) (0.0600,0.1921) (0.0833,0.2217) (0.1000,0.2562) (0.1200,0.2857) (0.1467,0.3153) (0.1567,0.3547) (0.1600,0.3547) (0.1633,0.3744) (0.1733,0.3744) (0.2100,0.3990) (0.2267,0.4039) (0.2400,0.4039) (0.2467,0.4039) (0.3233,0.5271) (0.3467,0.5616) (0.3667,0.6010) (0.3700,0.6158) (0.3767,0.6256) (0.4167,0.6453) (0.4200,0.6502) (0.4333,0.6502) (0.4367,0.6700) (0.4433,0.6700) (0.4500,0.6798) (0.4500,0.6897) (0.4600,0.6946) (0.4800,0.7044) (0.5167,0.7241) (0.5233,0.7340) (0.5467,0.7537) (0.5567,0.7586) (0.5867,0.7833) (0.6133,0.8128) (0.6133,0.8177) (0.6533,0.8325) (0.6733,0.8374) (0.7033,0.8621) (0.7200,0.8966) (1.0000,1.0000)};
\addlegendentry{DT (AUC 0.657)}
\addplot[thick, green!55!black] coordinates {(0.0000,0.0000) (0.0000,0.0049) (0.0000,0.0148) (0.0000,0.0345) (0.0033,0.0345) (0.0033,0.0443) (0.0067,0.0443) (0.0067,0.0640) (0.0100,0.0640) (0.0100,0.0788) (0.0133,0.0788) (0.0133,0.0985) (0.0167,0.0985) (0.0167,0.1084) (0.0200,0.1133) (0.0200,0.1330) (0.0233,0.1330) (0.0233,0.1527) (0.0267,0.1576) (0.0267,0.1626) (0.0300,0.1724) (0.0333,0.1724) (0.0333,0.1921) (0.0367,0.1921) (0.0367,0.2020) (0.0400,0.2020) (0.0400,0.2217) (0.0400,0.2266) (0.0400,0.2365) (0.0433,0.2463) (0.0433,0.2512) (0.0467,0.2562) (0.0467,0.2660) (0.0500,0.2660) (0.0500,0.2709) (0.0567,0.2759) (0.0600,0.2759) (0.0600,0.2906) (0.0633,0.3005) (0.0667,0.3005) (0.0700,0.3103) (0.0700,0.3202) (0.0733,0.3251) (0.0733,0.3350) (0.0733,0.3547) (0.0767,0.3645) (0.0800,0.3695) (0.0800,0.3744) (0.0867,0.3744) (0.0933,0.3793) (0.0967,0.3990) (0.1033,0.4187) (0.1067,0.4236) (0.1133,0.4236) (0.1133,0.4286) (0.1167,0.4286) (0.1200,0.4335) (0.1233,0.4433) (0.1267,0.4433) (0.1333,0.4433) (0.1400,0.4532) (0.1467,0.4581) (0.1533,0.4581) (0.1567,0.4581) (0.1567,0.4729) (0.1600,0.4729) (0.1733,0.4729) (0.1800,0.4778) (0.1800,0.4828) (0.1800,0.4926) (0.1800,0.4975) (0.2000,0.4975) (0.2033,0.5123) (0.2067,0.5320) (0.2067,0.5419) (0.2067,0.5567) (0.2100,0.5616) (0.2167,0.5616) (0.2333,0.5616) (0.2367,0.5665) (0.2467,0.5714) (0.2567,0.5714) (0.2700,0.5714) (0.2733,0.5764) (0.2733,0.5862) (0.2733,0.5911) (0.2767,0.5911) (0.2833,0.5961) (0.2833,0.6010) (0.2900,0.6010) (0.2933,0.6010) (0.2933,0.6059) (0.3000,0.6158) (0.3133,0.6207) (0.3167,0.6207) (0.3233,0.6404) (0.3367,0.6552) (0.3433,0.6650) (0.3533,0.6946) (0.3567,0.6995) (0.3700,0.6995) (0.3700,0.7044) (0.3867,0.7094) (0.3967,0.7094) (0.4033,0.7094) (0.4067,0.7094) (0.4133,0.7291) (0.4300,0.7291) (0.4300,0.7389) (0.4367,0.7488) (0.4467,0.7586) (0.4500,0.7586) (0.4633,0.7635) (0.4700,0.7734) (0.4800,0.7783) (0.4833,0.7833) (0.4933,0.7882) (0.5000,0.7931) (0.5067,0.7931) (0.5200,0.7931) (0.5233,0.7980) (0.5333,0.8030) (0.5433,0.8030) (0.5433,0.8177) (0.5433,0.8227) (0.5533,0.8227) (0.5600,0.8227) (0.5767,0.8276) (0.5900,0.8276) (0.6033,0.8276) (0.6233,0.8374) (0.6267,0.8424) (0.6467,0.8473) (0.6500,0.8473) (0.6567,0.8621) (0.6600,0.8621) (0.6667,0.8621) (0.6667,0.8719) (0.6700,0.8818) (0.6800,0.8818) (0.6800,0.8867) (0.6867,0.8916) (0.6933,0.8916) (0.6967,0.8966) (0.7000,0.8966) (0.7100,0.9113) (0.7233,0.9113) (0.7267,0.9113) (0.7267,0.9212) (0.7267,0.9261) (0.7333,0.9261) (0.7433,0.9310) (0.7467,0.9310) (0.7600,0.9310) (0.7633,0.9409) (0.7733,0.9458) (0.7767,0.9458) (0.7767,0.9557) (0.7833,0.9606) (0.7867,0.9606) (0.7933,0.9606) (0.8067,0.9655) (0.8133,0.9704) (0.8233,0.9704) (0.8300,0.9704) (0.8367,0.9754) (0.8500,0.9754) (0.8533,0.9754) (0.8600,0.9754) (0.8800,0.9754) (0.8833,0.9754) (0.8933,0.9754) (0.8967,0.9803) (0.9033,0.9852) (0.9067,0.9852) (0.9133,0.9852) (0.9167,0.9901) (0.9233,0.9901) (0.9367,0.9901) (0.9367,0.9951) (0.9467,0.9951) (0.9533,0.9951) (0.9633,0.9951) (0.9767,0.9951) (0.9800,1.0000) (1.0000,1.0000)};
\addlegendentry{RF (AUC 0.727)}
\addplot[thick, red!75!black] coordinates {(0.0000,0.0000) (0.0000,0.0049) (0.0000,0.0197) (0.0033,0.0197) (0.0033,0.0345) (0.0067,0.0345) (0.0067,0.0640) (0.0100,0.0640) (0.0100,0.1330) (0.0133,0.1330) (0.0133,0.1478) (0.0167,0.1478) (0.0167,0.1527) (0.0267,0.1527) (0.0267,0.1626) (0.0400,0.1626) (0.0400,0.1724) (0.0467,0.1724) (0.0467,0.2069) (0.0500,0.2069) (0.0500,0.2217) (0.0600,0.2217) (0.0600,0.2266) (0.0667,0.2266) (0.0667,0.2315) (0.0700,0.2315) (0.0700,0.2414) (0.0833,0.2414) (0.0833,0.2611) (0.0867,0.2611) (0.0867,0.2709) (0.0900,0.2709) (0.0900,0.2759) (0.0933,0.2759) (0.0933,0.2906) (0.0967,0.2906) (0.0967,0.3005) (0.1000,0.3005) (0.1000,0.3054) (0.1100,0.3054) (0.1100,0.3153) (0.1133,0.3153) (0.1133,0.3251) (0.1267,0.3251) (0.1267,0.3300) (0.1300,0.3300) (0.1300,0.3350) (0.1333,0.3350) (0.1333,0.3399) (0.1367,0.3399) (0.1367,0.3448) (0.1433,0.3448) (0.1433,0.3498) (0.1533,0.3498) (0.1533,0.3596) (0.1567,0.3596) (0.1567,0.3645) (0.1633,0.3645) (0.1633,0.3793) (0.1667,0.3793) (0.1667,0.3842) (0.1700,0.3842) (0.1700,0.3892) (0.1767,0.3892) (0.1767,0.3990) (0.1800,0.3990) (0.1800,0.4039) (0.2000,0.4039) (0.2000,0.4089) (0.2100,0.4089) (0.2100,0.4138) (0.2167,0.4138) (0.2167,0.4236) (0.2200,0.4236) (0.2200,0.4286) (0.2333,0.4286) (0.2333,0.4335) (0.2500,0.4335) (0.2500,0.4433) (0.2600,0.4433) (0.2600,0.4483) (0.2633,0.4483) (0.2633,0.4581) (0.2667,0.4581) (0.2667,0.4631) (0.2767,0.4631) (0.2767,0.4680) (0.2833,0.4680) (0.2833,0.4729) (0.3000,0.4729) (0.3000,0.4778) (0.3333,0.4778) (0.3333,0.4926) (0.3367,0.4926) (0.3367,0.5074) (0.3400,0.5074) (0.3400,0.5222) (0.3433,0.5222) (0.3433,0.5271) (0.3467,0.5271) (0.3467,0.5320) (0.3500,0.5320) (0.3500,0.5419) (0.3600,0.5419) (0.3600,0.5468) (0.3633,0.5468) (0.3633,0.5517) (0.3733,0.5517) (0.3733,0.5665) (0.3767,0.5665) (0.3767,0.5813) (0.3800,0.5813) (0.3800,0.5862) (0.3933,0.5862) (0.3933,0.6010) (0.4000,0.6010) (0.4000,0.6158) (0.4033,0.6158) (0.4033,0.6355) (0.4067,0.6355) (0.4067,0.6453) (0.4100,0.6453) (0.4100,0.6552) (0.4167,0.6552) (0.4167,0.6601) (0.4300,0.6601) (0.4300,0.6700) (0.4433,0.6700) (0.4433,0.6847) (0.4500,0.6847) (0.4500,0.6897) (0.4567,0.6897) (0.4567,0.7094) (0.4600,0.7094) (0.4600,0.7241) (0.4667,0.7241) (0.4667,0.7291) (0.4800,0.7291) (0.4800,0.7340) (0.4833,0.7340) (0.4833,0.7389) (0.5000,0.7389) (0.5000,0.7438) (0.5133,0.7438) (0.5133,0.7586) (0.5200,0.7586) (0.5200,0.7635) (0.5267,0.7635) (0.5267,0.7685) (0.5300,0.7685) (0.5300,0.7783) (0.5433,0.7783) (0.5433,0.7833) (0.5500,0.7833) (0.5500,0.7882) (0.5600,0.7882) (0.5600,0.8030) (0.5633,0.8030) (0.5633,0.8079) (0.5767,0.8079) (0.5767,0.8128) (0.5900,0.8128) (0.5900,0.8177) (0.5967,0.8177) (0.5967,0.8227) (0.6000,0.8227) (0.6000,0.8325) (0.6100,0.8325) (0.6100,0.8424) (0.6133,0.8424) (0.6133,0.8473) (0.6333,0.8473) (0.6333,0.8522) (0.6400,0.8522) (0.6400,0.8571) (0.6533,0.8571) (0.6533,0.8621) (0.6800,0.8621) (0.6800,0.8670) (0.6900,0.8670) (0.6900,0.8818) (0.7067,0.8818) (0.7067,0.8867) (0.7100,0.8867) (0.7100,0.8916) (0.7167,0.8916) (0.7167,0.8966) (0.7200,0.8966) (0.7200,0.9015) (0.7333,0.9015) (0.7333,0.9113) (0.7633,0.9113) (0.7633,0.9163) (0.7700,0.9163) (0.7700,0.9261) (0.7833,0.9261) (0.7833,0.9310) (0.8200,0.9310) (0.8200,0.9360) (0.8400,0.9360) (0.8400,0.9458) (0.8467,0.9458) (0.8467,0.9507) (0.8533,0.9507) (0.8533,0.9557) (0.8667,0.9557) (0.8667,0.9606) (0.8767,0.9606) (0.8767,0.9704) (0.8900,0.9704) (0.8900,0.9754) (0.9067,0.9754) (0.9067,0.9803) (0.9267,0.9803) (0.9267,0.9901) (0.9567,0.9901) (0.9567,0.9951) (0.9633,0.9951) (0.9633,1.0000) (1.0000,1.0000)};
\addlegendentry{XGBoost (AUC 0.671)}
\addplot[thick, violet!80!black] coordinates {(0.0000,0.0000) (0.0000,0.0148) (0.0133,0.0542) (0.0133,0.0591) (0.0233,0.0640) (0.0267,0.0690) (0.0667,0.1478) (0.1267,0.2709) (0.1400,0.2906) (0.1400,0.2956) (0.1433,0.3005) (0.1467,0.3103) (0.3300,0.4975) (0.3333,0.4975) (0.3733,0.5123) (0.4000,0.5616) (0.4067,0.5616) (0.4067,0.5665) (0.4100,0.5665) (0.4100,0.5714) (0.4133,0.5714) (0.9133,0.9803) (0.9233,0.9803) (0.9433,0.9803) (0.9433,0.9951) (0.9467,0.9951) (0.9667,0.9951) (0.9800,0.9951) (1.0000,1.0000)};
\addlegendentry{AdaBoost (AUC 0.615)}
\end{axis}\end{tikzpicture}
\caption{Potabilidad del agua: curvas ROC en test de los cinco modelos elegidos; clase positiva: Potable (1). La diagonal corresponde a un clasificador aleatorio.}
\label{fig:roc-potabilidad}
\end{figure}

\begin{figure}[H]\centering
\begin{tikzpicture}
\begin{axis}[ybar, bar width=11pt, width=0.9\linewidth, height=5.8cm, ymin=0, ymax=1.15, ytick={0,0.2,...,1},
 ylabel={F1 macro (validación cruzada)}, symbolic x coords={KNN,DT,RF,XGBoost,AdaBoost}, xtick=data,
 enlarge x limits=0.14, ymajorgrids, legend style={at={(0.5,-0.18)},anchor=north,legend columns=2,font=\small},
 nodes near coords, nodes near coords style={font=\tiny, rotate=90, anchor=west, /pgf/number format/fixed, /pgf/number format/precision=3}]
\addplot[fill=blue!50, draw=blue!70!black] coordinates {(KNN,0.6036) (DT,0.5908) (RF,0.6435) (XGBoost,0.6392) (AdaBoost,0.5404)};
\addplot[fill=teal!45, draw=teal!70!black] coordinates {(KNN,0.5364) (DT,0.5487) (RF,0.5436) (XGBoost,0.5439) (AdaBoost,0.5281)};
\legend{Todos los atributos (9),Selección por información mutua (5)}
\end{axis}\end{tikzpicture}
\caption{Potabilidad del agua: efecto de la selección de atributos sobre la mejor configuración de cada familia, medido con los mismos cinco pliegues.}
\label{fig:seleccion-potabilidad}
\end{figure}

\begin{figure}[H]\centering
\begin{tikzpicture}
\begin{axis}[xbar, width=0.8\linewidth, height=7.1cm, bar width=7pt,
 xlabel={Caída media del F1 macro en test al permutar}, ytick={0,1,2,3,4,5,6,7,8},
 yticklabels={{Conductivity},{Turbidity},{Organic\_\allowbreak{}carbon},{Trihalomethanes},{Solids},{Hardness},{Chloramines},{ph},{Sulfate}}, xmajorgrids, tick label style={font=\small}, enlarge y limits=0.08,
 scaled x ticks=false, xticklabel style={/pgf/number format/fixed, /pgf/number format/precision=2}]
\addplot[fill=green!45!black!60, draw=green!40!black, error bars/.cd, x dir=both, x explicit] coordinates {(0.0018,0) +- (0.0097,0) (0.0067,1) +- (0.0081,0) (0.0087,2) +- (0.0094,0) (0.0209,3) +- (0.0034,0) (0.0377,4) +- (0.0065,0) (0.0409,5) +- (0.0101,0) (0.0468,6) +- (0.0084,0) (0.0959,7) +- (0.0137,0) (0.1311,8) +- (0.0202,0)};
\end{axis}\end{tikzpicture}
\caption{Potabilidad del agua: importancia por permutación del modelo RF (media y desviación de 10 permutaciones de cada variable en test). Es una herramienta de interpretación: no participó en la selección.}
\label{fig:perm-potabilidad}
\end{figure}

Por validación cruzada, el mejor modelo es Random Forest con 500 árboles (F1 macro
$0.644\pm0.022$), y también es el mejor en prueba en F1 macro, exactitud y AUC. Aun
así, el problema es difícil: las AUC van de 0.615 a 0.727 y ninguna familia detecta
más del 50\,\% de las muestras potables. Hay 37 muestras potables de prueba que los
cinco modelos clasifican como no potables. Con medias casi idénticas entre clases, la
información útil está en regiones extremas y en combinaciones de variables, y los
modelos que mejor la aprovechan son los ensambles de árboles profundos.

La selección de atributos por información mutua empeoró las cinco familias, con caídas
de 0.01 (AdaBoost) a 0.10 (RF) en F1 macro (Figura~\ref{fig:seleccion-potabilidad}). El motivo es
instructivo. La información mutua calculada en train asigna 0.026 al sulfato, 0.014 a
la dureza y prácticamente 0 al \texttt{ph} y a las cloraminas, así que el filtro
descarta el \texttt{ph}. Sin embargo, la importancia por permutación muestra que el
\texttt{ph} es la segunda variable más importante para RF. Su efecto aparece sobre todo
en combinación con el sulfato (como muestran las reglas del árbol), y un filtro que
evalúa cada variable por separado no lo detecta. En este dataset conviene conservar
las nueve variables y dejar que los árboles encuentren las interacciones.

\FloatBarrier

\clearpage

\section{Resultados: Diabetes}

\subsection{Datos y partición}

Se analizan 390 pacientes con 17 variables: colesterol (\texttt{chol}), glucosa
estabilizada (\texttt{stab.glu}), HDL, razón colesterol/HDL (\texttt{ratio}), edad,
sexo, localidad, talla, peso, complexión (\texttt{frame}), dos mediciones de presión
arterial, cintura, cadera y tiempo desde la última comida (\texttt{time.ppn}). Sólo el
15.4\,\% son diabéticos (60 casos): en prueba hay 15 diabéticos, por lo que cada
diabético mal clasificado cambia el recall en 6.7 puntos, y en cada pliegue de
validación hay unos 9, lo que explica desviaciones de 0.06 a 0.12 en el F1 macro. La
glucosa estabilizada separa mucho mejor las clases que cualquier variable de
Potabilidad: su mediana es 186\,mg/dL en diabéticos y 86.5\,mg/dL en no diabéticos.
Tras la codificación one-hot hay 21 atributos; la información mutua de la glucosa
(0.228) es tres veces la del siguiente, \texttt{ratio} (0.063).

\begin{table}[H]\centering\small
\caption{Procedencia de Diabetes.}
\begin{tabular}{lp{0.62\linewidth}}\toprule
Elemento & Valor \\ \midrule
Archivo analizado & \texttt{diabetes.csv} \\
Original & \texttt{diabetes\_\allowbreak{}original.csv}, 403 filas \\
Filas analizadas & 390 \\
Objetivo & \texttt{diabetes} \\
Predictoras & 17 \\
Fuente & \url{https://hbiostat.org/data/repo/diabetes} \\
SHA-256 analizado & {\footnotesize\texttt{47b48720dd246d95528d251c31c27a76}\newline\texttt{f44c958a24faf1273b873658f4979e85}} \\
\bottomrule\end{tabular}
\end{table}

\begin{table}[H]\centering\small
\caption{Distribución de clases en Diabetes (partición estratificada 75/25, semilla 42).}
\begin{tabular}{lrrr}\toprule
Clase & Train & Test & Total \\ \midrule
No diabetes & 247 & 83 & 330 \\
Diabetes & 45 & 15 & 60 \\
\textbf{Total} & 292 & 98 & 390 \\
\bottomrule\end{tabular}
\end{table}

Una referencia trivial que siempre predice la clase mayoritaria de train (No diabetes) obtendría en test una exactitud de 84.7\% y un F1 macro de 0.459. Cualquier modelo útil debe superar con claridad ese F1 macro, no sólo la exactitud.

\subsection{K vecinos más cercanos (KNN)}

\begin{table}[H]\centering\footnotesize
\caption{K vecinos más cercanos (KNN) en Diabetes: validación cruzada estratificada de 5 pliegues sobre train. $\star$ = configuración elegida; $\Delta$ = diferencia de F1 macro frente al estándar.}
\begin{tabular}{>{\raggedright\arraybackslash}p{0.21\linewidth}>{\raggedright\arraybackslash}p{0.26\linewidth}rrrr}\toprule
Configuración & Cambio respecto del estándar & F1 macro CV & $\Delta$ & Exact. CV & Recall CV \\ \midrule
\texttt{estandar} & valores predeterminados de la biblioteca & 0.7190 $\pm$ 0.100 & -- & 89.0\% & 0.6808 \\
\texttt{k\_\allowbreak{}15} & $k=15$ vecinos (estándar: 5) & 0.6814 $\pm$ 0.099 & -0.038 & 88.7\% & 0.6424 \\
\texttt{pesos\_\allowbreak{}distancia} & votos ponderados por $1/d$ (estándar: uniformes) & 0.7190 $\pm$ 0.100 & +0.000 & 89.0\% & 0.6808 \\
\texttt{manhattan} & distancia Manhattan, $p=1$ (estándar: euclidiana) & 0.6358 $\pm$ 0.080 & -0.083 & 87.7\% & 0.6091 \\
\texttt{estandar+\allowbreak{}seleccion\_\allowbreak{}k11} $\star$ & mejor anterior + 11 atributos por información mutua & 0.7974 $\pm$ 0.070 & +0.078 & 91.8\% & 0.7424 \\
\bottomrule\end{tabular}
\end{table}

\textbf{Modelo final.} Configuración \texttt{estandar} con selección de atributos, reajustada con las 292 filas de train. F1 macro: 0.8547 en entrenamiento, 0.7974 en validación cruzada y 0.7478 en test (exactitud 88.8\%, AUC 0.754).

\begin{table}[H]\footnotesize\begin{minipage}[t]{0.44\linewidth}\centering\captionof{table}{K vecinos más cercanos (KNN): matriz de confusión en test (filas: real; columnas: predicción).}\begin{tabular}{lrr}\toprule
Real / pred. & Diabetes & No diabetes \\ \midrule
Diabetes & 7 & 8 \\
No diabetes & 3 & 80 \\
\bottomrule\end{tabular}\end{minipage}\hfill\begin{minipage}[t]{0.52\linewidth}\centering\captionof{table}{K vecinos más cercanos (KNN): métricas por clase en test.}\begin{tabular}{lrrrr}\toprule
Clase & Soporte & Precisión & Recall & F1 \\ \midrule
Diabetes & 15 & 0.7000 & 0.4667 & 0.5600 \\
No diabetes & 83 & 0.9091 & 0.9639 & 0.9357 \\
\bottomrule\end{tabular}\end{minipage}\end{table}

\paragraph{Análisis.} Con todos los atributos, KNN es la familia más débil (F1 macro 0.719 en el estándar).
Aumentar a $k=15$ empeora (0.681): con sólo unos 36 diabéticos en cada ajuste, en un
vecindario de 15 suele haber mayoría de no diabéticos y el voto favorece la clase
mayoritaria. La distancia Manhattan también empeora (0.636) y ponderar por distancia dio
exactamente las mismas métricas que el estándar en los cinco pliegues. El cambio que sí
importa es la selección de atributos: con los 11 de mayor información mutua el F1 macro
sube a 0.797 (+0.078), la mayor mejora de este dataset. Al eliminar variables poco
informativas como la talla, la localidad, el sexo o las segundas mediciones de presión
(ausentes en el 65\,\% de los pacientes e imputadas con la mediana), la distancia pasa
a reflejar sobre todo la glucosa. En prueba detecta 7 de 15 diabéticos (recall 0.467)
con precisión 0.70. Su AUC (0.754) es la más baja porque, con 5 vecinos, la
probabilidad sólo puede tomar seis valores (0, 0.2, \dots, 1), lo que produce una curva
ROC escalonada.

\FloatBarrier

\subsection{Árbol de decisión (DT)}

\begin{table}[H]\centering\footnotesize
\caption{Árbol de decisión (DT) en Diabetes: validación cruzada estratificada de 5 pliegues sobre train. $\star$ = configuración elegida; $\Delta$ = diferencia de F1 macro frente al estándar.}
\begin{tabular}{>{\raggedright\arraybackslash}p{0.21\linewidth}>{\raggedright\arraybackslash}p{0.26\linewidth}rrrr}\toprule
Configuración & Cambio respecto del estándar & F1 macro CV & $\Delta$ & Exact. CV & Recall CV \\ \midrule
\texttt{estandar} & valores predeterminados de la biblioteca & 0.7826 $\pm$ 0.065 & -- & 88.0\% & 0.8110 \\
\texttt{profundidad\_\allowbreak{}5} & profundidad máxima 5 (estándar: sin límite) & 0.7986 $\pm$ 0.117 & +0.016 & 89.4\% & 0.8101 \\
\texttt{hoja\_\allowbreak{}minima\_\allowbreak{}10} & mínimo 10 observaciones por hoja (estándar: 1) & 0.8064 $\pm$ 0.070 & +0.024 & 90.1\% & 0.8232 \\
\texttt{clases\_\allowbreak{}balanceadas} & pesos por observación inversos a la frecuencia de su clase & 0.7510 $\pm$ 0.077 & -0.032 & 87.7\% & 0.7454 \\
\texttt{hoja\_\allowbreak{}minima\_\allowbreak{}10+\allowbreak{}seleccion\_\allowbreak{}k11} $\star$ & mejor anterior + 11 atributos por información mutua & 0.8160 $\pm$ 0.085 & +0.033 & 90.8\% & 0.8182 \\
\bottomrule\end{tabular}
\end{table}

\textbf{Modelo final.} Configuración \texttt{hoja\_\allowbreak{}minima\_\allowbreak{}10} con selección de atributos, reajustada con las 292 filas de train. F1 macro: 0.8894 en entrenamiento, 0.8160 en validación cruzada y 0.7783 en test (exactitud 89.8\%, AUC 0.907).

\begin{table}[H]\footnotesize\begin{minipage}[t]{0.44\linewidth}\centering\captionof{table}{Árbol de decisión (DT): matriz de confusión en test (filas: real; columnas: predicción).}\begin{tabular}{lrr}\toprule
Real / pred. & Diabetes & No diabetes \\ \midrule
Diabetes & 8 & 7 \\
No diabetes & 3 & 80 \\
\bottomrule\end{tabular}\end{minipage}\hfill\begin{minipage}[t]{0.52\linewidth}\centering\captionof{table}{Árbol de decisión (DT): métricas por clase en test.}\begin{tabular}{lrrrr}\toprule
Clase & Soporte & Precisión & Recall & F1 \\ \midrule
Diabetes & 15 & 0.7273 & 0.5333 & 0.6154 \\
No diabetes & 83 & 0.9195 & 0.9639 & 0.9412 \\
\bottomrule\end{tabular}\end{minipage}\end{table}

\begin{figure}[H]\centering\begin{minipage}{0.92\linewidth}\scriptsize\begin{verbatim}
|--- stab.glu <= 158.00
|   |--- stab.glu <= 107.50
|   |   |--- age <= 63.50
|   |   |   |--- clase: No diabetes [n=183] (rama truncada)
|   |   |--- age > 63.50
|   |   |   |--- clase: No diabetes [n=30] (rama truncada)
|   |--- stab.glu > 107.50
|   |   |--- chol <= 243.50
|   |   |   |--- clase: No diabetes [n=33] (rama truncada)
|   |   |--- chol > 243.50
|   |   |   |--- clase: Diabetes [n=11]
|--- stab.glu > 158.00
|   |--- stab.glu <= 201.50
|   |   |--- clase: Diabetes [n=13]
|   |--- stab.glu > 201.50
|   |   |--- ratio <= 5.45
|   |   |   |--- clase: Diabetes [n=10]
|   |   |--- ratio > 5.45
|   |   |   |--- clase: Diabetes [n=12]
\end{verbatim}\end{minipage}\caption{Primeros tres niveles del árbol final (profundidad total 5, 11 hojas). Los umbrales están en las unidades originales: la estandarización se invirtió con la media y desviación de train. $n$ es la cantidad de observaciones de train en el nodo.}\end{figure}

\paragraph{Análisis.} Exigir al menos 10 observaciones por hoja es de nuevo la mejor variante (0.806, +0.024),
igual que en Potabilidad. Limitar la profundidad a 5 da un resultado parecido (0.799)
pero con una variación entre pliegues mucho mayor (0.117 frente a 0.070). Ponderar las clases empeora (0.751): con
tan pocos diabéticos, darles más peso hace que el árbol cree hojas para casos
aislados. La selección de atributos mejora un poco (0.816) y es la versión final. El
árbol resultante es pequeño (profundidad 5, 11 hojas) y fácil de leer: su primera regla
es \texttt{stab.glu} $>158$\,mg/dL $\Rightarrow$ diabetes, y en el rango de 107.5 a
158\,mg/dL marca como diabéticos a los pacientes con colesterol mayor que 243.5\,mg/dL. La importancia por
permutación confirma que en prueba el árbol depende exclusivamente de la glucosa (caída
de 0.252 en F1 macro; el resto, cero). En prueba detecta 8 de 15 diabéticos (recall
0.533, precisión 0.727) con exactitud 89.8\,\% y AUC 0.907. El umbral de 158\,mg/dL lo
aprendió el árbol de estos datos; no es un criterio diagnóstico. Como referencia, el
criterio clínico para una glucosa casual es 200\,mg/dL con síntomas \cite{ada}.

\FloatBarrier

\subsection{Random Forest (RF)}

\begin{table}[H]\centering\footnotesize
\caption{Random Forest (RF) en Diabetes: validación cruzada estratificada de 5 pliegues sobre train. $\star$ = configuración elegida; $\Delta$ = diferencia de F1 macro frente al estándar.}
\begin{tabular}{>{\raggedright\arraybackslash}p{0.21\linewidth}>{\raggedright\arraybackslash}p{0.26\linewidth}rrrr}\toprule
Configuración & Cambio respecto del estándar & F1 macro CV & $\Delta$ & Exact. CV & Recall CV \\ \midrule
\texttt{estandar} & valores predeterminados de la biblioteca & 0.8364 $\pm$ 0.087 & -- & 92.5\% & 0.8101 \\
\texttt{500\_\allowbreak{}arboles} & 500 árboles (estándar: 100) & 0.8177 $\pm$ 0.079 & -0.019 & 91.8\% & 0.7878 \\
\texttt{hoja\_\allowbreak{}minima\_\allowbreak{}5} & mínimo 5 observaciones por hoja (estándar: 1) & 0.7633 $\pm$ 0.068 & -0.073 & 90.1\% & 0.7232 \\
\texttt{clases\_\allowbreak{}balanceadas} & pesos por observación inversos a la frecuencia de su clase & 0.8162 $\pm$ 0.081 & -0.020 & 91.1\% & 0.8111 \\
\texttt{estandar+\allowbreak{}seleccion\_\allowbreak{}k11} $\star$ & mejor anterior + 11 atributos por información mutua & 0.8398 $\pm$ 0.066 & +0.003 & 92.5\% & 0.8192 \\
\bottomrule\end{tabular}
\end{table}

\textbf{Modelo final.} Configuración \texttt{estandar} con selección de atributos, reajustada con las 292 filas de train. F1 macro: 1.0000 en entrenamiento, 0.8398 en validación cruzada y 0.7638 en test (exactitud 88.8\%, AUC 0.892).

\begin{table}[H]\footnotesize\begin{minipage}[t]{0.44\linewidth}\centering\captionof{table}{Random Forest (RF): matriz de confusión en test (filas: real; columnas: predicción).}\begin{tabular}{lrr}\toprule
Real / pred. & Diabetes & No diabetes \\ \midrule
Diabetes & 8 & 7 \\
No diabetes & 4 & 79 \\
\bottomrule\end{tabular}\end{minipage}\hfill\begin{minipage}[t]{0.52\linewidth}\centering\captionof{table}{Random Forest (RF): métricas por clase en test.}\begin{tabular}{lrrrr}\toprule
Clase & Soporte & Precisión & Recall & F1 \\ \midrule
Diabetes & 15 & 0.6667 & 0.5333 & 0.5926 \\
No diabetes & 83 & 0.9186 & 0.9518 & 0.9349 \\
\bottomrule\end{tabular}\end{minipage}\end{table}

\paragraph{Análisis.} En validación cruzada, Random Forest es el mejor algoritmo del dataset. Su
configuración estándar (0.836) supera a las tres variantes: 500 árboles no mejora
(0.818), la ponderación de clases tampoco (0.816) y exigir 5 observaciones por hoja
empeora claramente (0.763), porque impide aislar a los pocos diabéticos en hojas
propias. Con selección de atributos alcanza 0.840 y una desviación menor (0.066 frente
a 0.087), así que esa es la versión final. En prueba tiene exactitud 88.8\,\%, F1
macro 0.764 y AUC 0.892; detecta 8 de 15 diabéticos y tiene un falso positivo más que
DT, XGBoost y AdaBoost (4 en lugar de 3). Por eso su F1 de prueba es ligeramente menor,
aunque fue el elegido por validación. Esa diferencia de un paciente está muy por debajo
de la incertidumbre del experimento. La importancia interna muestra que, además de la
glucosa (0.454), el bosque reparte peso entre \texttt{ratio}, HDL y edad, pero en
prueba sólo la glucosa produce una caída grande al permutarse.

\FloatBarrier

\subsection{XGBoost}

\begin{table}[H]\centering\footnotesize
\caption{XGBoost en Diabetes: validación cruzada estratificada de 5 pliegues sobre train. $\star$ = configuración elegida; $\Delta$ = diferencia de F1 macro frente al estándar.}
\begin{tabular}{>{\raggedright\arraybackslash}p{0.21\linewidth}>{\raggedright\arraybackslash}p{0.26\linewidth}rrrr}\toprule
Configuración & Cambio respecto del estándar & F1 macro CV & $\Delta$ & Exact. CV & Recall CV \\ \midrule
\texttt{estandar} & valores predeterminados de la biblioteca & 0.7921 $\pm$ 0.089 & -- & 90.1\% & 0.7869 \\
\texttt{tasa\_\allowbreak{}0\_\allowbreak{}05} $\star$ & tasa de aprendizaje 0.05 (estándar: 0.3) & 0.8208 $\pm$ 0.117 & +0.029 & 91.8\% & 0.8060 \\
\texttt{profundidad\_\allowbreak{}3} & profundidad máxima 3 (estándar: 6) & 0.8041 $\pm$ 0.098 & +0.012 & 90.8\% & 0.8000 \\
\texttt{clases\_\allowbreak{}balanceadas} & pesos por observación inversos a la frecuencia de su clase & 0.7885 $\pm$ 0.102 & -0.004 & 89.4\% & 0.8010 \\
\texttt{tasa\_\allowbreak{}0\_\allowbreak{}05+\allowbreak{}seleccion\_\allowbreak{}k11} & mejor anterior + 11 atributos por información mutua & 0.8034 $\pm$ 0.108 & +0.011 & 91.1\% & 0.7929 \\
\bottomrule\end{tabular}
\end{table}

\textbf{Modelo final.} Configuración \texttt{tasa\_\allowbreak{}0\_\allowbreak{}05} con todos los atributos, reajustada con las 292 filas de train. F1 macro: 0.9934 en entrenamiento, 0.8208 en validación cruzada y 0.7783 en test (exactitud 89.8\%, AUC 0.929).

\begin{table}[H]\footnotesize\begin{minipage}[t]{0.44\linewidth}\centering\captionof{table}{XGBoost: matriz de confusión en test (filas: real; columnas: predicción).}\begin{tabular}{lrr}\toprule
Real / pred. & Diabetes & No diabetes \\ \midrule
Diabetes & 8 & 7 \\
No diabetes & 3 & 80 \\
\bottomrule\end{tabular}\end{minipage}\hfill\begin{minipage}[t]{0.52\linewidth}\centering\captionof{table}{XGBoost: métricas por clase en test.}\begin{tabular}{lrrrr}\toprule
Clase & Soporte & Precisión & Recall & F1 \\ \midrule
Diabetes & 15 & 0.7273 & 0.5333 & 0.6154 \\
No diabetes & 83 & 0.9195 & 0.9639 & 0.9412 \\
\bottomrule\end{tabular}\end{minipage}\end{table}

\paragraph{Análisis.} Bajar la tasa de aprendizaje a 0.05 es la mejor variante (0.821, +0.029): con sólo 292
pacientes de entrenamiento, que cada árbol aporte menos funciona como regularización.
Reducir la profundidad a 3 ayuda algo menos (0.804) y ponderar las clases no ayuda
(0.789). A diferencia de las otras familias, la selección de atributos empeora a XGBoost
(0.803): el algoritmo ya elige en cada división la variable más útil y quitarle
variables secundarias le resta poco ruido y algo de información. En prueba obtiene la
mejor AUC del estudio (0.929), exactitud 89.8\,\% y F1 macro 0.778, con la misma matriz
de confusión que DT y AdaBoost. Su F1 de entrenamiento (0.993) indica que aun con tasa
baja casi memoriza el entrenamiento.

\FloatBarrier

\subsection{AdaBoost}

\begin{table}[H]\centering\footnotesize
\caption{AdaBoost en Diabetes: validación cruzada estratificada de 5 pliegues sobre train. $\star$ = configuración elegida; $\Delta$ = diferencia de F1 macro frente al estándar.}
\begin{tabular}{>{\raggedright\arraybackslash}p{0.21\linewidth}>{\raggedright\arraybackslash}p{0.26\linewidth}rrrr}\toprule
Configuración & Cambio respecto del estándar & F1 macro CV & $\Delta$ & Exact. CV & Recall CV \\ \midrule
\texttt{estandar} & valores predeterminados de la biblioteca & 0.7920 $\pm$ 0.058 & -- & 90.4\% & 0.7616 \\
\texttt{200\_\allowbreak{}estimadores} & 200 estimadores (estándar: 50) & 0.7896 $\pm$ 0.079 & -0.002 & 90.4\% & 0.7616 \\
\texttt{tasa\_\allowbreak{}0\_\allowbreak{}5} & tasa de aprendizaje 0.5 (estándar: 1.0) & 0.7994 $\pm$ 0.122 & +0.007 & 91.1\% & 0.7838 \\
\texttt{clases\_\allowbreak{}balanceadas} & pesos por observación inversos a la frecuencia de su clase & 0.7944 $\pm$ 0.063 & +0.002 & 89.4\% & 0.8009 \\
\texttt{tasa\_\allowbreak{}0\_\allowbreak{}5+\allowbreak{}seleccion\_\allowbreak{}k11} $\star$ & mejor anterior + 11 atributos por información mutua & 0.8027 $\pm$ 0.087 & +0.011 & 91.1\% & 0.7747 \\
\bottomrule\end{tabular}
\end{table}

\textbf{Modelo final.} Configuración \texttt{tasa\_\allowbreak{}0\_\allowbreak{}5} con selección de atributos, reajustada con las 292 filas de train. F1 macro: 0.8865 en entrenamiento, 0.8027 en validación cruzada y 0.7783 en test (exactitud 89.8\%, AUC 0.922).

\begin{table}[H]\footnotesize\begin{minipage}[t]{0.44\linewidth}\centering\captionof{table}{AdaBoost: matriz de confusión en test (filas: real; columnas: predicción).}\begin{tabular}{lrr}\toprule
Real / pred. & Diabetes & No diabetes \\ \midrule
Diabetes & 8 & 7 \\
No diabetes & 3 & 80 \\
\bottomrule\end{tabular}\end{minipage}\hfill\begin{minipage}[t]{0.52\linewidth}\centering\captionof{table}{AdaBoost: métricas por clase en test.}\begin{tabular}{lrrrr}\toprule
Clase & Soporte & Precisión & Recall & F1 \\ \midrule
Diabetes & 15 & 0.7273 & 0.5333 & 0.6154 \\
No diabetes & 83 & 0.9195 & 0.9639 & 0.9412 \\
\bottomrule\end{tabular}\end{minipage}\end{table}

\paragraph{Análisis.} Las variantes de AdaBoost quedan todas entre 0.790 y 0.799, sin diferencias claras. La
tasa 0.5 es la de mayor media (0.799), pero también la de mayor desviación (0.122). La
ponderación de clases logra el mayor recall macro de validación (0.801), aunque no el
mayor F1. Con selección de atributos sube a 0.803 y esa es la versión final. En prueba
coincide con DT y XGBoost (exactitud 89.8\,\%, F1 macro 0.778, 8 de 15 diabéticos) y
tiene AUC 0.922. Es el ensamble con menor sobreajuste (F1 0.887 en entrenamiento frente
a 0.778 en prueba). En este dataset los stumps sí funcionan, a diferencia de
Potabilidad: un único corte en la glucosa ya separa bien las clases, y los siguientes
stumps afinan con el colesterol (importancia interna 0.120) y el peso (0.106).

\FloatBarrier

\subsection{Comparación de los cinco modelos elegidos}

\begin{table}[H]\centering\footnotesize
\caption{Diabetes: resumen de los ganadores. La selección usó sólo la columna F1 CV; el mejor algoritmo por validación es RF.}
\begin{tabular}{llrrrrrr}\toprule
Familia & Configuración & F1 train & F1 CV & F1 test & Exactitud test & Recall Diabetes & AUC \\ \midrule
KNN & \texttt{estandar} + sel. & 0.8547 & 0.7974 & 0.7478 & 88.8\% & 0.4667 & 0.754 \\
DT & \texttt{hoja\_\allowbreak{}minima\_\allowbreak{}10} + sel. & 0.8894 & 0.8160 & 0.7783 & 89.8\% & 0.5333 & 0.907 \\
RF & \texttt{estandar} + sel. & 1.0000 & 0.8398 & 0.7638 & 88.8\% & 0.5333 & 0.892 \\
XGBoost & \texttt{tasa\_\allowbreak{}0\_\allowbreak{}05} & 0.9934 & 0.8208 & 0.7783 & 89.8\% & 0.5333 & 0.929 \\
AdaBoost & \texttt{tasa\_\allowbreak{}0\_\allowbreak{}5} + sel. & 0.8865 & 0.8027 & 0.7783 & 89.8\% & 0.5333 & 0.922 \\
\bottomrule\end{tabular}
\end{table}

\begin{figure}[H]\centering
\begin{tikzpicture}
\begin{axis}[ybar, bar width=8pt, width=\linewidth, height=6.2cm, ymin=0, ymax=1.08,
 ylabel={F1 macro}, symbolic x coords={KNN,DT,RF,XGBoost,AdaBoost}, xtick=data,
 enlarge x limits=0.14, ymajorgrids, legend style={at={(0.5,-0.16)},anchor=north,legend columns=4,font=\small},
 tick label style={font=\small}]
\addplot[fill=gray!35, draw=gray!70!black] coordinates {(KNN,0.8547) (DT,0.8894) (RF,1.0000) (XGBoost,0.9934) (AdaBoost,0.8865)};
\addplot[fill=blue!55, draw=blue!70!black, error bars/.cd, y dir=both, y explicit] coordinates {(KNN,0.7974) +- (0,0.0704) (DT,0.8160) +- (0,0.0846) (RF,0.8398) +- (0,0.0659) (XGBoost,0.8208) +- (0,0.1170) (AdaBoost,0.8027) +- (0,0.0869)};
\addplot[fill=orange!70, draw=orange!80!black] coordinates {(KNN,0.7478) (DT,0.7783) (RF,0.7638) (XGBoost,0.7783) (AdaBoost,0.7783)};
\draw[dashed, red!70!black, thick] ({rel axis cs:0,0}|-{axis cs:KNN,0.4586}) -- ({rel axis cs:1,0}|-{axis cs:KNN,0.4586});
\legend{Entrenamiento, Validación cruzada (media $\pm$ d.e.), Test}
\end{axis}\end{tikzpicture}
\caption{Diabetes: F1 macro del modelo elegido en cada familia. La línea roja es la referencia que siempre predice la clase mayoritaria (F1 macro 0.459). La distancia entre la barra gris y la azul indica sobreajuste.}
\label{fig:resumen-diabetes}
\end{figure}

\begin{figure}[H]\centering
\begin{tikzpicture}
\begin{axis}[width=0.72\linewidth, height=0.62\linewidth, xmin=0, xmax=1, ymin=0, ymax=1,
 xlabel={Tasa de falsos positivos}, ylabel={Tasa de verdaderos positivos}, grid=major,
 legend style={at={(0.98,0.02)}, anchor=south east, font=\footnotesize}]
\addplot[dashed, gray, forget plot] coordinates {(0,0) (1,1)};
\addplot[thick, blue!70!black] coordinates {(0.0000,0.0000) (0.0000,0.0667) (0.0241,0.2667) (0.0361,0.4667) (0.0482,0.5333) (0.2771,0.6667) (1.0000,1.0000)};
\addlegendentry{KNN (AUC 0.754)}
\addplot[thick, orange!85!black] coordinates {(0.0000,0.0000) (0.0000,0.0667) (0.0120,0.4000) (0.0361,0.5333) (0.0361,0.8667) (0.1325,0.8667) (1.0000,1.0000)};
\addlegendentry{DT (AUC 0.907)}
\addplot[thick, green!55!black] coordinates {(0.0000,0.0000) (0.0120,0.0000) (0.0241,0.0667) (0.0241,0.2000) (0.0361,0.2000) (0.0361,0.5333) (0.0602,0.5333) (0.0602,0.6000) (0.1084,0.6000) (0.1084,0.7333) (0.1566,0.7333) (0.1566,0.8667) (0.1687,0.8667) (0.1928,0.8667) (0.2289,0.9333) (0.2651,0.9333) (0.2892,0.9333) (0.3855,0.9333) (0.4096,0.9333) (0.4578,0.9333) (0.5301,0.9333) (0.6145,1.0000) (0.8313,1.0000) (1.0000,1.0000)};
\addlegendentry{RF (AUC 0.892)}
\addplot[thick, red!75!black] coordinates {(0.0000,0.0000) (0.0000,0.0667) (0.0000,0.4667) (0.0361,0.4667) (0.0361,0.8000) (0.0482,0.8000) (0.0482,0.8667) (0.0964,0.8667) (0.0964,0.9333) (0.7470,0.9333) (0.7470,1.0000) (0.8193,1.0000) (0.8434,1.0000) (1.0000,1.0000)};
\addlegendentry{XGBoost (AUC 0.929)}
\addplot[thick, violet!80!black] coordinates {(0.0000,0.0000) (0.0000,0.0667) (0.0000,0.1333) (0.0120,0.1333) (0.0120,0.2000) (0.0241,0.2000) (0.0241,0.3333) (0.0361,0.3333) (0.0361,0.6000) (0.0361,0.7333) (0.0482,0.7333) (0.0482,0.8667) (0.0482,0.9333) (0.1325,0.9333) (0.1687,0.9333) (0.1928,0.9333) (0.2169,0.9333) (0.2410,0.9333) (0.2651,0.9333) (0.3012,0.9333) (0.3253,0.9333) (0.4096,0.9333) (0.4578,0.9333) (0.5181,0.9333) (0.5663,0.9333) (0.5904,0.9333) (0.6386,0.9333) (0.6627,0.9333) (0.7229,0.9333) (0.7470,0.9333) (0.7470,1.0000) (0.8434,1.0000) (0.9036,1.0000) (0.9398,1.0000) (0.9639,1.0000) (1.0000,1.0000)};
\addlegendentry{AdaBoost (AUC 0.922)}
\end{axis}\end{tikzpicture}
\caption{Diabetes: curvas ROC en test de los cinco modelos elegidos; clase positiva: Diabetes. La diagonal corresponde a un clasificador aleatorio.}
\label{fig:roc-diabetes}
\end{figure}

\begin{figure}[H]\centering
\begin{tikzpicture}
\begin{axis}[ybar, bar width=11pt, width=0.9\linewidth, height=5.8cm, ymin=0, ymax=1.15, ytick={0,0.2,...,1},
 ylabel={F1 macro (validación cruzada)}, symbolic x coords={KNN,DT,RF,XGBoost,AdaBoost}, xtick=data,
 enlarge x limits=0.14, ymajorgrids, legend style={at={(0.5,-0.18)},anchor=north,legend columns=2,font=\small},
 nodes near coords, nodes near coords style={font=\tiny, rotate=90, anchor=west, /pgf/number format/fixed, /pgf/number format/precision=3}]
\addplot[fill=blue!50, draw=blue!70!black] coordinates {(KNN,0.7190) (DT,0.8064) (RF,0.8364) (XGBoost,0.8208) (AdaBoost,0.7994)};
\addplot[fill=teal!45, draw=teal!70!black] coordinates {(KNN,0.7974) (DT,0.8160) (RF,0.8398) (XGBoost,0.8034) (AdaBoost,0.8027)};
\legend{Todos los atributos (21),Selección por información mutua (11)}
\end{axis}\end{tikzpicture}
\caption{Diabetes: efecto de la selección de atributos sobre la mejor configuración de cada familia, medido con los mismos cinco pliegues.}
\label{fig:seleccion-diabetes}
\end{figure}

\begin{figure}[H]\centering
\begin{tikzpicture}
\begin{axis}[xbar, width=0.8\linewidth, height=7.1cm, bar width=7pt,
 xlabel={Caída media del F1 macro en test al permutar}, ytick={0,1,2,3,4,5,6,7,8},
 yticklabels={{frame},{height},{gender},{location},{age},{bp.1s},{weight},{bp.1d},{stab.glu}}, xmajorgrids, tick label style={font=\small}, enlarge y limits=0.08,
 scaled x ticks=false, xticklabel style={/pgf/number format/fixed, /pgf/number format/precision=2}]
\addplot[fill=green!45!black!60, draw=green!40!black, error bars/.cd, x dir=both, x explicit] coordinates {(0.0000,0) +- (0.0000,0) (0.0000,1) +- (0.0000,0) (0.0000,2) +- (0.0000,0) (0.0000,3) +- (0.0000,0) (0.0028,4) +- (0.0058,0) (0.0040,5) +- (0.0094,0) (0.0070,6) +- (0.0139,0) (0.0241,7) +- (0.0168,0) (0.2487,8) +- (0.0692,0)};
\end{axis}\end{tikzpicture}
\caption{Diabetes: importancia por permutación del modelo RF (media y desviación de 10 permutaciones de cada variable en test). Es una herramienta de interpretación: no participó en la selección.}
\label{fig:perm-diabetes}
\end{figure}

Los cuatro métodos basados en árboles quedan entre 0.80 y 0.84 de F1 macro en
validación cruzada, con desviaciones de 0.07 a 0.12. Las diferencias entre ellos son
menores que esa variación, así que el orden no es concluyente; RF se eligió por tener
la mayor media. En prueba todos superan claramente la referencia (F1 macro 0.459) y DT,
XGBoost y AdaBoost producen la misma matriz de confusión: 8 diabéticos detectados, 7 no
detectados y 3 falsos positivos.

Los 7 diabéticos que no detectan DT, RF, XGBoost y AdaBoost son exactamente los mismos
(KNN falla esos 7 y uno más), y todos tienen glucosa estabilizada entre 90 y 130\,mg/dL, mientras que los 8 detectados
tienen 196\,mg/dL o más. La etiqueta se definió con la hemoglobina glicosilada, que
refleja la glucosa promedio de los últimos dos o tres meses, y puede estar elevada
aunque una medición puntual de glucosa sea normal. Como \texttt{glyhb} se excluyó para
evitar fuga, esos pacientes no se distinguen de los no diabéticos con las variables
disponibles. Es un límite de la información, no de los algoritmos: ningún ajuste de
hiperparámetros los recuperaría sin aumentar mucho los falsos positivos.

Las AUC de 0.89 a 0.93 indican que los modelos ordenan bien a los pacientes por riesgo.
Bajar el umbral de decisión de 0.5 aumentaría el recall de diabetes a cambio de más
falsos positivos. Ese umbral debería elegirse con validación cruzada y no con el
conjunto de prueba, por lo que se deja como trabajo futuro. A diferencia de
Potabilidad, aquí la selección de atributos ayudó a cuatro de cinco familias
(Figura~\ref{fig:seleccion-diabetes}): la información relevante se concentra en una
variable con relación monótona con la clase, y la información mutua la detecta sin
problema.

\FloatBarrier

\clearpage

\section{Discusión comparativa}

\subsection{Qué variantes ganaron y por qué}
La Tabla~\ref{tab:variantes} resume la configuración elegida en cada caso. Hay dos
patrones que se repiten en ambos datasets. Para el árbol de decisión, exigir un mínimo
de 10 observaciones por hoja fue siempre la mejor variante: un árbol sin límites
memoriza el entrenamiento, y limitar el tamaño de las hojas reduce esa varianza sin
eliminar las reglas útiles. Para KNN, aumentar $k$ a 15 no fue la mejor variante en
ningún dataset: en Diabetes bajó el F1 macro de 0.719 a 0.681, porque con 45 diabéticos
en train un vecindario de 15 casos queda casi siempre dominado por la clase mayoritaria;
en Potabilidad subió apenas 0.011, menos que el cambio a distancia Manhattan.

En cambio, otras decisiones dependen del dataset. Ponderar las clases fue la mejor
variante de XGBoost y AdaBoost en Potabilidad, donde los modelos tienden a predecir
``no potable''. En Diabetes no fue la mejor variante en ninguna familia y bajó el F1 de
DT (0.783 a 0.751) y RF (0.836 a 0.816): dar más peso a 45 diabéticos lleva a los
árboles a crear reglas para casos aislados. La selección de atributos mejoró cuatro
de cinco familias en Diabetes (en RF sólo 0.003, una diferencia despreciable) y empeoró las cinco en Potabilidad. La explicación es la
forma de la relación entre variables y clase: en Diabetes la señal está casi toda en la
glucosa y es monótona (más glucosa, más riesgo), así que un filtro univariado la
encuentra; en Potabilidad la señal está en los extremos y en la combinación de sulfato
y \texttt{ph}, y el filtro descarta el \texttt{ph} aunque sea la segunda variable más
importante para el mejor modelo.

\begin{table}[H]\centering\small
\caption{Configuración elegida por validación cruzada en cada familia y dataset.}
\label{tab:variantes}
\begin{tabular}{lll}\toprule
Familia & Potabilidad & Diabetes \\ \midrule
KNN & Manhattan ($p=1$) & estándar + selección de 11 atributos \\
DT & mínimo 10 por hoja & mínimo 10 por hoja + selección \\
RF & 500 árboles & estándar + selección \\
XGBoost & clases balanceadas & tasa de aprendizaje 0.05 \\
AdaBoost & clases balanceadas & tasa 0.5 + selección \\ \bottomrule
\end{tabular}
\end{table}

\subsection{Sobreajuste y estabilidad}
Las Figuras~\ref{fig:resumen-potabilidad} y~\ref{fig:resumen-diabetes} comparan el F1 de
entrenamiento, validación cruzada y prueba. Random Forest y XGBoost llegan a F1 de
entrenamiento de 1.0 o casi 1.0 en ambos datasets: memorizan el entrenamiento. Aun así,
RF es el que mejor generaliza en Potabilidad, porque el promedio de cientos de árboles
cancela buena parte de sus errores individuales. AdaBoost con stumps es el extremo
opuesto en Potabilidad: casi no sobreajusta (0.575 en entrenamiento y 0.582 en prueba),
pero subajusta. KNN y DT quedan en medio.

La validación cruzada también muestra cuánto se puede confiar en cada diferencia. En
Potabilidad (unos 300 casos por pliegue) las desviaciones de los ganadores están entre
0.02 y 0.05, y la ventaja de RF sobre KNN, DT y AdaBoost es clara; frente a XGBoost
(0.639) la diferencia es menor que la desviación. En Diabetes (unos 58 casos y 9
diabéticos por pliegue) las desviaciones llegan a 0.12 y las cuatro familias de árboles
son estadísticamente indistinguibles. Una sola partición de validación, como la que se
usaba antes en el framework, habría dado un ganador sin advertir esa incertidumbre.

\subsection{Exactitud frente a F1 macro}
En Diabetes la exactitud de prueba (88.8--89.8\,\%) parece alta, pero la referencia que
siempre predice ``no diabetes'' ya obtiene 84.7\,\%. El F1 macro (0.75--0.78 frente a
0.459) y el recall de diabetes (0.47--0.53) muestran mejor lo que hace el modelo: detecta
la mitad de los diabéticos. En Potabilidad, AdaBoost estándar tiene una exactitud de
validación de 59.7\,\%, igual a la proporción de la clase mayoritaria: sin el F1 macro
se habría pasado por alto que el modelo casi no aprendió nada.

\subsection{Limitaciones}
\begin{itemize}
\item Se usó una única partición de prueba (503 y 98 casos). Con 15 diabéticos en
prueba, un solo paciente cambia el recall en 6.7 puntos. No se calcularon intervalos de
confianza ni pruebas de significancia.
\item El catálogo tiene tres variantes por familia y cada una cambia un solo
hiperparámetro. Es deliberado, para poder atribuir los efectos, pero no explora
combinaciones (por ejemplo, más árboles y hojas mínimas a la vez).
\item Todos los modelos usan el umbral de probabilidad 0.5. Ajustarlo con validación
cruzada podría aumentar el recall de la clase minoritaria.
\item La clase de Diabetes se construyó con el criterio histórico $\text{glyhb}>7$, y la
versión de Potabilidad usada son los casos completos del archivo público. Los resultados
son educativos: no constituyen un diagnóstico médico ni una certificación de potabilidad.
\item La importancia por permutación mide de qué depende el modelo, no relaciones
causales.
\end{itemize}


\section{Conclusiones}

\begin{enumerate}
\item Los cinco algoritmos quedaron integrados en la clase \texttt{Clasificacion} del
framework, con un protocolo reproducible: preparación aprendida sólo con
entrenamiento, validación cruzada estratificada de 5 pliegues para elegir y una única
evaluación en prueba. Cada experimento se repite con un comando y deja toda la evidencia.

\item En Potabilidad, el mejor modelo es Random Forest con 500 árboles (F1 macro de
validación $0.644\pm0.022$; en prueba, F1 macro 0.659, exactitud 70.2\,\% y AUC 0.727).
Supera con claridad la referencia (59.6\,\% de exactitud) y a las demás familias, pero
detecta sólo el 43\,\% de las muestras potables. Las variables apenas difieren en
promedio entre clases y la señal aparece en los extremos y en la combinación de sulfato
y \texttt{ph}, algo que los ensambles de árboles profundos captan mejor que KNN o
AdaBoost con stumps.

\item En Diabetes, DT, RF, XGBoost y AdaBoost obtienen resultados equivalentes dentro de
la variación de la validación cruzada (F1 macro 0.80--0.84). RF fue el elegido por
validación; en prueba los cuatro detectan 8 de 15 diabéticos con AUC entre 0.89 y 0.93.
La glucosa estabilizada explica casi todo el poder predictivo, y los siete diabéticos
que ningún modelo detecta tienen glucosa normal en la medición, lo que marca un límite
de la información disponible.

\item Las variantes que mejoraron dependen del algoritmo y del dataset. Exigir hojas de al menos
10 observaciones mejoró el árbol de decisión en ambos datasets, pero limitar las hojas
de Random Forest empeoró el bosque, porque el promedio ya controla la varianza. Ponderar
las clases ayudó a los métodos de boosting cuando tendían a ignorar la clase minoritaria
(Potabilidad). Seleccionar atributos por información mutua ayudó
cuando la señal estaba en pocas variables con relación monótona (Diabetes) y perjudicó
cuando dependía de interacciones (Potabilidad).

\item La exactitud sola es engañosa con clases desbalanceadas. El F1 macro, el recall
por clase y la comparación con una referencia trivial fueron necesarios para interpretar
los resultados.
\end{enumerate}

Como trabajo futuro se propone elegir el umbral de decisión con validación cruzada para
aumentar el recall de la clase minoritaria, explorar combinaciones de hiperparámetros con
búsqueda en rejilla dentro de la validación cruzada y repetir el experimento con varias
semillas para estimar la estabilidad de las conclusiones.


\clearpage

\begin{thebibliography}{99}
\bibitem{hastie}
T. Hastie, R. Tibshirani y J. Friedman. \emph{The Elements of Statistical Learning:
Data Mining, Inference, and Prediction}, 2.\textsuperscript{a} ed. Springer, 2009.

\bibitem{coverhart}
T. M. Cover y P. E. Hart. Nearest Neighbor Pattern Classification.
\emph{IEEE Transactions on Information Theory}, 13(1):21--27, 1967.
\url{https://doi.org/10.1109/TIT.1967.1053964}.

\bibitem{sklearnknn}
scikit-learn. \emph{Nearest Neighbors}, documentación oficial.
\url{https://scikit-learn.org/stable/modules/neighbors.html}.

\bibitem{sklearntree}
scikit-learn. \emph{Decision Trees}, documentación oficial.
\url{https://scikit-learn.org/stable/modules/tree.html}.

\bibitem{breiman}
L. Breiman. Random Forests. \emph{Machine Learning}, 45:5--32, 2001.
\url{https://doi.org/10.1023/A:1010933404324}.

\bibitem{sklearnensemble}
scikit-learn. \emph{Ensembles: Gradient boosting, random forests, bagging, voting,
stacking}, documentación oficial.
\url{https://scikit-learn.org/stable/modules/ensemble.html}.

\bibitem{chen}
T. Chen y C. Guestrin. XGBoost: A Scalable Tree Boosting System. En
\emph{Proceedings of the 22nd ACM SIGKDD International Conference on Knowledge
Discovery and Data Mining}, pp. 785--794, 2016.
\url{https://doi.org/10.1145/2939672.2939785}.

\bibitem{xgboostmodel}
XGBoost. \emph{Introduction to Boosted Trees}, documentación oficial.
\url{https://xgboost.readthedocs.io/en/stable/tutorials/model.html}.

\bibitem{freund}
Y. Freund y R. E. Schapire. A Decision-Theoretic Generalization of On-Line Learning
and an Application to Boosting. \emph{Journal of Computer and System Sciences},
55(1):119--139, 1997. \url{https://doi.org/10.1006/jcss.1997.1504}.

\bibitem{sklearnada}
scikit-learn. \emph{AdaBoostClassifier}, documentación oficial.
\url{https://scikit-learn.org/stable/modules/generated/sklearn.ensemble.AdaBoostClassifier.html}.

\bibitem{sklearnmetrics}
scikit-learn. \emph{Metrics and scoring: quantifying the quality of predictions},
documentación oficial.
\url{https://scikit-learn.org/stable/modules/model_evaluation.html}.

\bibitem{sklearnleakage}
scikit-learn. \emph{Common pitfalls and recommended practices: Data leakage},
documentación oficial.
\url{https://scikit-learn.org/stable/common_pitfalls.html}.

\bibitem{kohavi}
R. Kohavi. A Study of Cross-Validation and Bootstrap for Accuracy Estimation and
Model Selection. En \emph{Proceedings of the 14th International Joint Conference on
Artificial Intelligence (IJCAI)}, pp. 1137--1145, 1995.

\bibitem{sklearncv}
scikit-learn. \emph{Cross-validation: evaluating estimator performance},
documentación oficial.
\url{https://scikit-learn.org/stable/modules/cross_validation.html}.

\bibitem{kraskov}
A. Kraskov, H. Stögbauer y P. Grassberger. Estimating mutual information.
\emph{Physical Review E}, 69:066138, 2004.
\url{https://doi.org/10.1103/PhysRevE.69.066138}.

\bibitem{ross}
B. C. Ross. Mutual Information between Discrete and Continuous Data Sets.
\emph{PLoS ONE}, 9(2):e87357, 2014.
\url{https://doi.org/10.1371/journal.pone.0087357}.

\bibitem{guyon}
I. Guyon y A. Elisseeff. An Introduction to Variable and Feature Selection.
\emph{Journal of Machine Learning Research}, 3:1157--1182, 2003.

\bibitem{hegarcia}
H. He y E. A. Garcia. Learning from Imbalanced Data. \emph{IEEE Transactions on
Knowledge and Data Engineering}, 21(9):1263--1284, 2009.
\url{https://doi.org/10.1109/TKDE.2008.239}.

\bibitem{fawcett}
T. Fawcett. An introduction to ROC analysis. \emph{Pattern Recognition Letters},
27(8):861--874, 2006. \url{https://doi.org/10.1016/j.patrec.2005.10.010}.

\bibitem{sklearnperm}
scikit-learn. \emph{Permutation feature importance}, documentación oficial.
\url{https://scikit-learn.org/stable/modules/permutation_importance.html}.

\bibitem{pedregosa}
F. Pedregosa et al. Scikit-learn: Machine Learning in Python. \emph{Journal of
Machine Learning Research}, 12:2825--2830, 2011.

\bibitem{datospotabilidad}
A. Kadiwal. \emph{Water Quality} (water\_potability.csv). Kaggle, 2021.
\url{https://www.kaggle.com/datasets/adityakadiwal/water-potability}.

\bibitem{datosdiabetes}
Vanderbilt University, Department of Biostatistics. \emph{Diabetes Dataset}
(J. P. Willems, J. T. Saunders, D. E. Hunt y J. B. Schorling).
\url{https://hbiostat.org/data/repo/diabetes}.

\bibitem{ada}
American Diabetes Association Professional Practice Committee. 2. Diagnosis and
Classification of Diabetes: Standards of Care in Diabetes---2024. \emph{Diabetes
Care}, 47(Suppl. 1):S20--S42, 2024. \url{https://doi.org/10.2337/dc24-S002}.
\end{thebibliography}


\clearpage

\appendix

\section{Reproducción de los experimentos}

Desde la carpeta raíz del proyecto, con Python 3.11 o superior y las dependencias de
\texttt{requirements-lab02-lock.txt} instaladas, los pasos son los siguientes. Cada
experimento tarda alrededor de un minuto. La carpeta de salida debe estar vacía: el
script no sobrescribe evidencia anterior.

\begin{verbatim}
python -m pip install -r requirements-lab02-lock.txt

# 1. Preparación de los datos (determinista; conserva los originales)
python -m scripts.preparar_potabilidad_lab02 data/lab02/water_potability.csv
python -m scripts.preparar_diabetes_lab02 data/lab02/diabetes_original.csv

# 2. Experimentos (validación cruzada de 5 pliegues, semilla 42)
python -m scripts.ejecutar_lab02 data/lab02/potabilidad.csv --target Potability
    --salida salida_lab02/potabilidad --semilla 42
    --fuente https://www.kaggle.com/datasets/adityakadiwal/water-potability
python -m scripts.ejecutar_lab02 data/lab02/diabetes.csv --target diabetes
    --features chol stab.glu hdl ratio location age gender height weight
    frame bp.1s bp.1d bp.2s bp.2d waist hip time.ppn --incluir-categoricas
    --salida salida_lab02/diabetes --semilla 42
    --fuente https://hbiostat.org/data/repo/diabetes

# 3. Informe y pruebas del software
python -m scripts.generar_informe_lab02 salida_lab02 --autor "Nombre (ID)"
python -m pytest -q
\end{verbatim}

Las líneas indentadas continúan el comando anterior (en PowerShell se escriben en una
sola línea). Cada carpeta de resultados contiene \texttt{configuracion.json} (argumentos,
SHA-256 del CSV, versiones, índices de train, test y de cada pliegue),
\texttt{comparacion\_validacion.csv} (todas las configuraciones con media, desviación y
puntuación por pliegue), \texttt{mejores\_modelos.json}, \texttt{ranking\_atributos.csv}
y, por algoritmo, la matriz de confusión, las predicciones de prueba, la curva ROC y la
importancia por permutación; además, \texttt{DT/arbol.txt} con las reglas del árbol.


\begin{table}[H]\centering\small
\caption{Versiones registradas durante la ejecución.}
\begin{tabular}{ll}\toprule
Paquete & Versión \\ \midrule
numpy & 2.5.3 \\
pandas & 3.0.5 \\
scikit-learn & 1.9.1 \\
xgboost & 3.4.1 \\
python & 3.13.16 \\
\bottomrule\end{tabular}
\end{table}

\begin{table}[H]\centering\small
\caption{Huellas de los datos utilizados.}
\begin{tabular}{lp{0.6\linewidth}}\toprule
Archivo & SHA-256 \\ \midrule
water\_\allowbreak{}potability.csv & {\footnotesize\texttt{904004bde729bfe3d2e195f46343bcee}\newline\texttt{ad09e32a0eb95bb8184e7e20e029b2bf}} \\
potabilidad.csv & {\footnotesize\texttt{d2569bece50025b9dbf21cfd720ecbed}\newline\texttt{cdaa838ca082db29ef853c5fb9da8e12}} \\
diabetes\_\allowbreak{}original.csv & {\footnotesize\texttt{3de140ba6771ee5af7043500065c667c}\newline\texttt{fb5c059eb840cc22e14cabf1ebfc0838}} \\
diabetes.csv & {\footnotesize\texttt{47b48720dd246d95528d251c31c27a76}\newline\texttt{f44c958a24faf1273b873658f4979e85}} \\
\bottomrule\end{tabular}
\end{table}

\end{document}
